\subsection{Weyl inequalities}

In this issue, we consider Hilbert spaces over K = C.

Let E be a Hilbert space. For every \(n \in N\) , we recall that the Hilbertian exterior power \(\hat{\Lambda}^{n}E\) was defined in EVT, V, p. 34. For every Hilbert space F and for \(u \in \mathcal{L}(E;F)\) , we have also defined the linear map \(\hat{\Lambda}^{n}u \in \mathcal{L}(\hat{\Lambda}^{n}E;\hat{\Lambda}^{n}F)\) (loc. cit.). These constructions are functorial: for every Hilbert space G and for every linear map \(v \in \mathcal{L}(F;G)\) , the formula

\[
\widehat {\Lambda} ^ {n} v \circ \widehat {\Lambda} ^ {n} u = \widehat {\Lambda} ^ {n} (v \circ u)
\]

is valid (loc. cit., formula (28)).

Let H be a closed subspace of E and let \(i_{H}\) be the canonical injection from H into E. For every endomorphism u of E, we denote by \(u_{H}\) the endomorphism \(i_{H}^{*}ui_{H}\) of H.

Lemma 2. --- Let \(n \in N\) . The map \(\hat{\Lambda}^{n}i_{H}\) is a linear isometric mapping from \(\hat{\Lambda}^{n}H\) into \(\hat{\Lambda}^{n}E\) .

Let \((e_{j})_{j\in\mathrm{J}}\) be an orthonormal basis of H and \((e_{j})_{j\in\mathrm{J}^{\prime}}\) an orthonormal basis of E, where \(J\subset J^{\prime}\) . Equip \(J^{\prime}\) with a total order. The elements \(e_{j_{1}}\wedge\cdots\wedge e_{j_{n}}\) for \(j_{1}<\cdots<j_{n}\) in \(J^{\prime}\) (resp. in J) form an orthonormal basis of \(\widehat{\Lambda}^{n}E\) (resp. of \(\widehat{\Lambda}^{n}H\) ) by prop. 5 of EVT, V, p. 34, prop. 5. The lemma follows.

In what follows, we shall identify \(\hat{\Lambda}^n\mathrm{H}\) with a closed subspace of \(\hat{\Lambda}^n\mathrm{E}\) by means of the mapping \(\hat{\Lambda}^n i_{\mathrm{H}}\) .

Lemma 3. --- Let F be a Hilbert space and \(u \in \mathcal{L}(\mathrm{E};\mathrm{F})\). Let \(n \in \mathbf{N}\) . a) We have ( \(\widehat{\Lambda}^n u\) ) \(^*\) = \(\widehat{\Lambda}^n (u^*)\) ;

\begin{enumerate}
\def\labelenumi{\alph{enumi})}
\setcounter{enumi}{1}
\item
  If F = E and u is Hermitian (resp. positive, normal, unitary), then \(\hat{\Lambda}^{n}u\) is Hermitian (resp. positive, normal, unitary);
\item
  We have \(|\widehat{\Lambda}^n u| = \widehat{\Lambda}^n |u|\) ;
\item
  Let H be a closed subspace of E. The restriction \((\widehat{\Lambda}^{n}u)|\widehat{\Lambda}^{n}\mathrm{H}\) of \(\widehat{\Lambda}^{n}u\) to \(\widehat{\Lambda}^{n}\mathrm{H}\) is equal to \(\widehat{\Lambda}^{n}(u|\mathrm{H})\) (equality in \(\mathcal{L}(\widehat{\Lambda}^{n}\mathrm{H};\widehat{\Lambda}^{n}\mathrm{F})\) );
\item
  Suppose F = E. Let H be a closed subspace of E. We have \((\hat{\Lambda}^{n}u)_{\hat{\Lambda}^{n}\mathrm{H}}=\hat{\Lambda}^{n}(u_{\mathrm{H}})\) in \(\mathcal{L}(\hat{\Lambda}^{n}\mathrm{H})\) .
\end{enumerate}

The assertion \(a\) ) follows from EVT, V, p. 39, formula (13). It immediately implies that \(\widehat{\Lambda}^{n}u\) is Hermitian (resp. unitary) when \(u\) is Hermitian (resp. unitary). If \(u\) is positive, then \(u^{1/2}\) is Hermitian, hence so is \(\widehat{\Lambda}^{n}(u^{1/2})\) ; the relation \(\widehat{\Lambda}^{n}u = (\widehat{\Lambda}^{n}u^{1/2})^{2}\) implies that \(\widehat{\Lambda}^{n}u\) is positive (I, p. 138, prop. 8). This proves \(b\) ).

What precedes allows one to compute that

\[
\begin{array}{r l} & {\left(\widehat {\Lambda} ^ {n} | u |\right) ^ {2} = \widehat {\Lambda} ^ {n} (| u | ^ {2}) = \widehat {\Lambda} ^ {n} (u ^ {*} u)} \\ & {\qquad = \widehat {\Lambda} ^ {n} u ^ {*} \widehat {\Lambda} ^ {n} u = \left(\widehat {\Lambda} ^ {n} u\right) ^ {*} (\widehat {\Lambda} ^ {n} u) = | \widehat {\Lambda} ^ {n} u | ^ {2}.} \end{array}
\]

Since \(\widehat{\Lambda}^n |u|\) is positive by \(b\) ), it follows from prop. 16 of I, p. 118 that \(\widehat{\Lambda}^n |u| = |\widehat{\Lambda}^n u|\) , whence \(c\) ).

Let \(i_{\mathrm{H}}\) be the canonical injection of H into E. The space \(\widehat{\Lambda}^{n}\mathrm{H}\) is identified with a closed subspace of \(\widehat{\Lambda}^{n}\mathrm{E}\) by the mapping \(\widehat{\Lambda}^{n}i_{\mathrm{H}}\) , whence

\[
(\widehat {\Lambda} ^ {n} u) | \widehat {\Lambda} ^ {n} \mathrm{H} = \widehat {\Lambda} ^ {n} u \circ \widehat {\Lambda} ^ {n} i _ {\mathrm{H}} = \widehat {\Lambda} ^ {n} (u \circ i _ {\mathrm{H}}) = \widehat {\Lambda} ^ {n} (u | \mathrm{H}),
\]

and

\[
(\widehat {\Lambda} ^ {n} u) _ {\widehat {\Lambda} ^ {n} \mathrm{H}} = (\widehat {\Lambda} ^ {n} i _ {\mathrm{H}}) ^ {*} \circ \widehat {\Lambda} ^ {n} u \circ \widehat {\Lambda} ^ {n} i _ {\mathrm{H}} = \widehat {\Lambda} ^ {n} (i _ {\mathrm{H}} ^ {*} u i _ {\mathrm{H}}) = \widehat {\Lambda} ^ {n} (u _ {\mathrm{H}}),
\]

which proves \(d\) ) and \(e\) ).

Let F be a Hilbert space, and let \(u \in \mathcal{L}^{c}(\mathrm{E};\mathrm{F})\) . As in number 3 of IV, p. 151, we denote by \(I_{E}\) the set of dimensions of finite-dimensional subspaces F of E such that \(F \neq E\) . We recall that \((\alpha_{n}(u))_{n \in I_{E}}\) denotes the extended sequence of singular values of u (def. 3 of IV, p. 156).

PROPOSITION 10. --- One has the equality

\[
\prod_ {i = 0} ^ {n} \alpha_ {i} (u) = \left\| \widehat {\Lambda} ^ {n + 1} u \right\|
\]

for all \(n \in I_{E}\) .

Lemma 3, c) shows that \(\|\widehat{\Lambda}^{n+1}u\|=\|\widehat{\Lambda}^{n+1}|u\|\) ; moreover, since \(\alpha_{i}(u)=\alpha_{i}(|u|)\) for all \(i\in I_{E}\) (remark 5 of IV, p. 157, (2)), it suffices to prove the assertion for \(|u|\) . This allows us to assume that F=E and that u is positive.

Let then \(\mathrm{B}=(e_{j})_{j\in\mathrm{J}}\) be an orthonormal basis of E such that u is diagonalizable in the basis B (theorem 1 of IV, p. 149), and let \((\lambda_{j})_{j\in\mathrm{J}}\) be the family of eigenvalues of u in the basis B. Endow J with a total order, and denote by \(J_{n}\) the set of strictly increasing families \((j_{0},\ldots,j_{n})\in\mathbf{J}^{n+1}\) of elements of J. The vectors

\[
e _ {\iota} = e _ {j _ {0}} \wedge \dots \wedge e _ {j _ {n}}
\]

for \(\iota=(j_{0},\ldots,j_{n})\in\mathrm{J}_{n}\) form an orthonormal basis \(B_{n}\) of \(\widehat{\Lambda}^{n+1}\mathrm{E}\) (EVT, V, p. 34, prop. 5). For every \(\iota=(j_{0},\ldots,j_{n})\in\mathrm{J}_{n}\) , denote

\[
\lambda_ {\iota} = \prod_ {j = 0} ^ {n} \lambda_ {i _ {j}}.
\]

It follows that \((\hat{\Lambda}^{n + 1}u)e_{\iota} = \lambda_{\iota}e_{\iota}\) , hence \(\hat{\Lambda}^{n + 1}u\) is diagonal in the basis \(\mathrm{B}_n\) . Consequently, one has

\[
\| \widehat {\Lambda} ^ {n + 1} u \| = \sup _ {\iota \in \mathrm{I} _ {n}} \lambda_ {\iota}
\]

(lemma 1 of IV, p. 147), which is equal to the product \(\lambda_{0}(u)\cdots\lambda_{n}(u)\) of the \(n+1\) largest eigenvalues of \(u\) . The desired formula follows since \(\lambda_{i}(u)=\alpha_{i}(u)\) for all \(i\in\mathrm{I}_{\mathrm{E}}\) when \(u\) is positive (remark 5 of IV, p. 157, (3)).

In particular, if \(\alpha_{1}(u) < \alpha_{0}(u) = \| u\|\) , one sees that the inequality \(\| \widehat{\Lambda}^{n}(u)\| \leqslant \| u\|^{n}\) (EVT, V, p. 34, formula (29)) is not an equality in general if \(n \geqslant 2\) .

COROLLARY. --- Let G be a Hilbert space and \(v \in \mathcal{L}^{\mathrm{c}}(\mathrm{F};\mathrm{G})\) . We have

\[
\prod_ {i = 0} ^ {n} \alpha_ {i} (v \circ u) \leqslant \prod_ {i = 0} ^ {n} \alpha_ {i} (u) \alpha_ {i} (v)
\]

for all \(n \in \mathrm{I}_{\mathrm{E}}\) .

It suffices to note that

\[
\| \widehat {\Lambda} ^ {n + 1} (v u) \| = \| \widehat {\Lambda} ^ {n + 1} v \circ \widehat {\Lambda} ^ {n + 1} u \| \leqslant \| \widehat {\Lambda} ^ {n + 1} v \| \| \widehat {\Lambda} ^ {n + 1} u \|
\]

and to apply proposition 10.

Lemma 4. --- Let A be a ring. Let \(n \in \mathbf{N}\) and \((a_i)_{0 \leqslant i \leqslant n}\) and \((b_i)_{0 \leqslant i \leqslant n}\) be families of elements of A. For \(0 \leqslant j \leqslant n\) set

\[
\mathrm{A} _ {j} = \sum_ {i = 0} ^ {j} a _ {i}.
\]

One has

\[
\sum_ {i = 0} ^ {n} a _ {i} b _ {i} = \mathrm{A} _ {n} b _ {n} - \sum_ {i = 0} ^ {n - 1} \mathrm{A} _ {i} (b _ {i + 1} - b _ {i}).
\]

Set \(A_{-1}=0\) . We have \(a_{i}=A_{i}-A_{i-1}\) for \(0\leqslant i\leqslant n\) , hence

\[
\sum_ {i = 0} ^ {n} a _ {i} b _ {i} = \sum_ {i = 0} ^ {n} (\mathrm{A} _ {i} - \mathrm{A} _ {i - 1}) b _ {i} = \sum_ {i = 0} ^ {n - 1} \mathrm{A} _ {i} (b _ {i} - b _ {i + 1}) + \mathrm{A} _ {n} b _ {n},
\]

as desired.

Let I be an interval of \(\overline{R}\) not containing \(+\infty\) . Recall (FVR, I, p. 38, remark) that a continuous function \(f: I \to [-\infty, +\infty[\) is said to be convex if its restriction to the interior of I is a convex function with values in R; it then has a limit (finite or infinite) on the right at inf I. In the statements below, the product of 0 and an element of \(\{-\infty, +\infty\}\) is by convention equal to 0.

Lemma 5. --- Let I ⊂ R be an interval not containing +∞ and f an increasing convex function from I into \([-\infty, +\infty[\).

Let n be a natural integer and

\[
a _ {0} \geqslant a _ {1} \geqslant \dots \geqslant a _ {n}, \quad b _ {0} \geqslant b _ {1} \geqslant \dots \geqslant b _ {n}
\]

elements of I.

Let \(\varrho_{0} \geqslant \varrho_{1} \geqslant \cdots \geqslant \varrho_{n}\) be positive real numbers. Suppose that

\[
\sum_ {i = 0} ^ {j} a _ {i} \leqslant \sum_ {i = 0} ^ {j} b _ {i}.\tag{5}
\]

for every integer j such that \(0 \leqslant j \leqslant n\) . We then have

\[
\sum_ {i = 0} ^ {n} \varrho_ {i} f (a _ {i}) \leqslant \sum_ {i = 0} ^ {n} \varrho_ {i} f (b _ {i}).\tag{6}
\]

Suppose first that \(a_i \in \mathring{\mathrm{I}}\) , and hence \(f(a_i) \in \mathbf{R}\) , for all \(i\). Let \(i\) be such that \(0 \leqslant i \leqslant n\) , and let \((\alpha_i, \beta_i)\) be real numbers such that the line with equation \(y = \alpha_i x + \beta_i\) is a supporting line of the graph of \(f\) at the point \((a_i, f(a_i))\) (FVR, I, p. 37). One consequently has \(f(a_i) = \alpha_i a_i + \beta_i\) and \(f(b_i) \geqslant \alpha_i b_i + \beta_i\) since the graph of \(f\) lies above the supporting line, whence

\[
f (b _ {i}) - f (a _ {i}) \geqslant \alpha_ {i} (b _ {i} - a _ {i}).\tag{7}
\]

Moreover, if \(i < n\) , we have

\[
\alpha_ {i} \geqslant f _ {g} ^ {\prime} (a _ {i}) \geqslant f _ {d} ^ {\prime} (a _ {i + 1}) \geqslant \alpha_ {i + 1}
\]

(loc. cit. and FVR, I, p. 36, cor. 1); furthermore \(\alpha_{i} \geqslant 0\) since \(f\) is increasing (FVR, I, p. 22, corollary), whence \(\varrho_{i}\alpha_{i} \geqslant \varrho_{i+1}\alpha_{i+1} \geqslant 0\) for \(0 \leqslant i < n\) .

Let \(j\) be an integer such that \(0 \leqslant j \leqslant n\) . Set

\[
\mathrm{A} _ {j} = \sum_ {i = 0} ^ {j} (b _ {i} - a _ {i}),
\]

so that \(A_{j} \geqslant 0\) by hypothesis (5). Applying inequality (7) and then Lemma 4, we deduce

\[
\begin{array}{l} \sum_ {i = 0} ^ {n} \varrho_ {i} (f (b _ {i}) - f (a _ {i})) \geqslant \sum_ {i = 0} ^ {n} \varrho_ {i} \alpha_ {i} (b _ {i} - a _ {i}) \\ \qquad = \varrho_ {n} \alpha_ {n} A _ {n} + \sum_ {j = 0} ^ {n - 1} (\varrho_ {j} \alpha_ {j} - \varrho_ {j + 1} \alpha_ {j + 1}) A _ {j} \geqslant 0. \end{array}
\]

Consider the general case, and reason by induction on \(n\). If one of the \(a_{i}\) does not belong to the interior of I, we necessarily have \(a_{0} = \sup I\) or \(a_{n} = \inf I\) .

Suppose first that \(a_{n}=\inf I\) . We then have \(a_{n}\leqslant b_{n}\) and therefore \(\varrho_{n}f(a_{n})\leqslant\varrho_{n}f(b_{n})\) . The induction hypothesis, applied to the families \((a_{0},\ldots,a_{n-1})\) , \((b_{0},\ldots,b_{n-1})\) and \((\varrho_{0},\ldots,\varrho_{n-1})\) , implies

\[
\sum_ {i = 0} ^ {n - 1} \varrho_ {i} f (a _ {i}) \leqslant \sum_ {i = 0} ^ {n - 1} \varrho_ {i} f (b _ {i}),
\]

whence the desired inequality follows by adding \(\varrho_{n}f(a_{n})\). The case where \(a_{0} = \sup I\) is treated in a similar manner.

PROPOSITION 11 (Weyl inequalities). --- Let G be a Hilbert space and \(v \in \mathcal{L}^{\mathrm{c}}(\mathrm{F};\mathrm{G})\). Let \(g: \mathbf{R}_{+} \to [-\infty, +\infty[\) be a nondecreasing function such that the function \(g \circ \exp\) is convex. We have

\[
\sum_ {i = 0} ^ {n} g \left(\alpha_ {i} (v \circ u)\right) \leqslant \sum_ {i = 0} ^ {n} g \left(\alpha_ {i} (v) \alpha_ {i} (u)\right)
\]

for all \(n \in I_{E} \cap I_{F}\) .

Set \(I = [-\infty, +\infty[\) and \(f = g \circ \exp\). It is permissible to apply Lemma 5 with

\[
a _ {i} = \log (\alpha_ {i} (v \circ u)) \in \mathrm{I}, \quad b _ {i} = \log (\alpha_ {i} (v) \alpha_ {i} (u)) \in \mathrm{I}
\]

and \(\varrho_{i} = 1\) for \(0\leqslant i\leqslant n\) since

\[
\sum_ {i = 0} ^ {j} a _ {i} = \log \Bigl (\prod_ {i = 0} ^ {j} \alpha_ {i} (v \circ u) \Bigr) \leqslant \log \Bigl (\prod_ {i = 0} ^ {j} \alpha_ {i} (v) \alpha_ {i} (u) \Bigr) = \sum_ {i = 0} ^ {j} b _ {i}
\]

for \(0 \leqslant j \leqslant n\) by the corollary of prop. 10. Inequality (6) is then the desired conclusion.

Lemma 6. --- Let g and h be convex functions defined on intervals I and J of R, respectively. If g is increasing and defined on the image of h, then the function g ∘ h is convex on J.

Indeed, for \(t \in [0,1]\) and \((x,y) \in \mathrm{J} \times \mathrm{J}\) , we have

\[
\begin{array}{c} g (h (t x + (1 - t) y)) \leqslant g (t h (x) + (1 - t) h (y)) \\ \leqslant t g (h (x)) + (1 - t) g (h (y)). \end{array}
\]

COROLLARY. --- Let G be a Hilbert space and \(v \in \mathcal{L}^{\mathrm{c}}(\mathrm{F};\mathrm{G})\). Let \(n \in I_{E} \cap I_{F}\) .

\begin{enumerate}
\def\labelenumi{\alph{enumi})}
\tightlist
\item
  Let \(r \in \mathbf{R}_{+}^{*}\) . We have
\end{enumerate}

\[
\sum_ {i = 0} ^ {n} \alpha_ {i} (v \circ u) ^ {r} \leqslant \sum_ {i = 0} ^ {n} \alpha_ {i} (v) ^ {r} \alpha_ {i} (u) ^ {r}.
\]

\begin{enumerate}
\def\labelenumi{\alph{enumi})}
\setcounter{enumi}{1}
\tightlist
\item
  Let \(p, q, r \in \mathbf{R}_{+}^{*}\) such that \(\frac{1}{p} + \frac{1}{q} = \frac{1}{r}\) . Then
\end{enumerate}

\[
\left(\sum_ {i = 0} ^ {n} \alpha_ {i} (v \circ u) ^ {r}\right) ^ {1 / r} \leqslant \left(\sum_ {i = 0} ^ {n} \alpha_ {i} (v) ^ {p}\right) ^ {1 / p} \left(\sum_ {i = 0} ^ {n} \alpha_ {i} (u) ^ {q}\right) ^ {1 / q}.
\]

\begin{enumerate}
\def\labelenumi{\alph{enumi})}
\setcounter{enumi}{2}
\tightlist
\item
  Suppose that F = E. For every integer \(m \geqslant 2\) , we have
\end{enumerate}

\[
\sum_ {i = 0} ^ {n} \alpha_ {i} (u ^ {m}) \leqslant \left(\sum_ {i = 0} ^ {n} \alpha_ {i} (u) ^ {2}\right) ^ {m / 2}.
\]

Let \(r \in R_{+}^{*}\) and let g be the function defined on \(\mathbf{R}_{+}\) by \(g(x) = x^{r}\) . The function \(g \circ \exp\) is convex (lemma 6), so assertion a) follows from prop. 11 applied to the function g.

The assertion \(b\) ) follows from \(a\) ) and from Hölder's inequality (INT, I, prop. 4).

Let \(m \geqslant 2\) be an integer. Let us apply \(b)\) with \(r = 1\) , \(p = q = 2\) and \(v = u^{m-1}\). We find

\[
\sum_ {i = 0} ^ {n} \alpha_ {i} (u ^ {m}) \leqslant \left(\sum_ {i = 0} ^ {n} \alpha_ {i} (u ^ {m - 1}) ^ {2}\right) ^ {1 / 2} \left(\sum_ {i = 0} ^ {n} \alpha_ {i} (u) ^ {2}\right) ^ {1 / 2}.
\]

Let us prove c) by induction on \(m \geqslant 2\) . The preceding inequality establishes the assertion when m = 2. Suppose that \(m \geqslant 3\) and that the assertion is valid for m - 1; since we have

\[
\sum_ {i = 0} ^ {n} \alpha_ {i} (u ^ {m - 1}) ^ {2} \leqslant \left(\sum_ {i = 0} ^ {m} \alpha_ {i} (u ^ {m - 1})\right) ^ {2},
\]

the above inequality implies

\[
\sum_ {i = 0} ^ {n} \alpha_ {i} (u ^ {m}) \leqslant \left(\sum_ {i = 0} ^ {n} \alpha_ {i} (u ^ {m - 1})\right) ^ {1 / 2} \left(\sum_ {i = 0} ^ {n} \alpha_ {i} (u) ^ {2}\right) ^ {1 / 2} \leqslant \left(\sum_ {i = 0} ^ {n} \alpha_ {i} (u) ^ {2}\right) ^ {m / 2}
\]

by applying the induction hypothesis.

\subsection{Endomorphisms of finite trace}

We recall that a positive endomorphism \(u\) of E is of finite trace if and only if there exists an orthonormal basis \((e_i)_{i \in I}\) of E such that

\[
\sum_ {i \in \mathrm{I}} \langle e _ {i} | u (e _ {i}) \rangle <   + \infty
\]

(EVT, V, p. 48, lemma 3 and p. 49, def. 7). If K = C, the space \(\mathcal{L}_1(\mathrm{E})\) of endomorphisms of finite trace of E is the vector space generated by the set of positive endomorphisms of finite trace (EVT, V, p. 50, def. 8); if K = R, the space \(\mathcal{L}_1(\mathrm{E})\) is defined as the intersection \(\mathcal{L}(\mathrm{E}) \cap \mathcal{L}_1(\mathrm{E}_{(\mathbf{C})})\) (EVT, V, p. 50).

If \(u\in \mathcal{L}_1(\mathrm{E})\) , then the series

\[
\sum_ {i \in \mathrm{I}} \langle e _ {i} | u (e _ {i}) \rangle
\]

converges for every orthonormal basis \((e_{i})_{i\in\mathrm{I}}\) of E and its sum is independent of the orthonormal basis; this sum is called the trace \(\operatorname{Tr}(u)\) of u (EVT, V, p. 50). If K = R, we have \(\operatorname{Tr}(u) = \operatorname{Tr}(u_{(\mathbf{C})})\) .

Let \(u \in \mathcal{L}_1(\mathrm{E})\) . We have \(u^* \in \mathcal{L}_1(\mathrm{E})\) and \(\operatorname{Tr}(u^*) = \overline{\operatorname{Tr}(u)}\) (loc. cit.).

PROPOSITION 12. --- Let \(B = (e_i)_{i \in I}\) be an orthonormal basis of E. Let \(u \in \mathcal{D}_B(E)\) and denote by \(\lambda = (\lambda_i)_{i \in I}\) the family of its eigenvalues. The endomorphism \(u\) is of finite trace if and only if the family \(\lambda\) is summable in K. We then have \(\operatorname{Tr}(u) = \sum \lambda_i\) .

Replacing \(u\) by \(u_{(\mathbf{C})}\) if necessary, we may assume that \(K = \mathbf{C}\) .

Suppose first that \(u\) is of finite trace. According to EVT, V, p. 50 and p. 49, formula (25), the family \((\langle e_i | u(e_i) \rangle)_{i \in I} = (\lambda_i)_{i \in I}\) is summable.

Conversely, suppose that the family \(\lambda\) is summable.

Each of the families \((\mathcal{R}(\lambda_{i})^{+})\) , \((\mathcal{R}(\lambda_{i})^{-})\) , \((\mathcal{I}(\lambda_{i})^{+})\) , \((\mathcal{I}(\lambda_{i})^{-})\) is then summable. The endomorphism u is a linear combination of the elements of \(\mathcal{D}_{\mathrm{B}}(\mathrm{E})\) whose eigenvalues are these families. These elements of \(\mathcal{D}_{\mathrm{B}}(\mathrm{E})\) are positive. Since, by definition, the space of finite-trace endomorphisms is generated by positive finite-trace endomorphisms, we may therefore assume that \(\lambda_{i} \geqslant 0\) for all i. Since \(\langle e_{i} | u(e_{i}) \rangle = \lambda_{i}\), the family \((\langle e_{i} | u(e_{i}) \rangle)_{i \in \mathrm{I}}\) is summable, hence u is of finite trace (EVT, V, p. 48, lemma 3).

Finally, if u has finite trace, then by EVT, V, p. 50, we have

\[
\operatorname{Tr} (u) = \sum_ {i \in \mathrm{I}} \left\langle e _ {i} \mid u \left(e _ {i}\right) \right\rangle = \sum_ {i \in \mathrm{I}} \lambda_ {i}.
\]

COROLLARY 1. --- Let u be an endomorphism of E with finite trace.

\begin{enumerate}
\def\labelenumi{\alph{enumi})}
\item
  The endomorphism u is compact;
\item
  Let \(\mathrm{B} = (e_i)_{i \in \mathrm{I}}\) be an orthonormal basis of the initial space \(\operatorname{Ker}(u)^{\circ}\) of u, \(\mathrm{C} = (f_i)_{i \in \mathrm{I}}\) be an orthonormal family of E and \((\alpha_i)_{i \in \mathrm{I}}\) be a family in \((\mathbf{R}_+^*)^{\mathrm{I}}\) such that \(u(e_i) = \alpha_i f_i\) for all \(i \in I\) (cor. 2 of IV, p. 150). We have
\end{enumerate}

\[
\operatorname{Tr} (u) = \sum_ {i \in I} \alpha_ {i} \langle e _ {i} | f _ {i} \rangle .
\]

To prove \(a\) ), one may assume that \(\mathbf{K} = \mathbf{C}\) (EVT, V, p. 50 and remark 4 of III, p. 2) and that \(u\) is positive (EVT, V, p. 50, def. 8).

By EVT, V, p. 56, cor. 1, there exists an orthonormal basis \(B = (e_i)_{i \in I}\) of E such that \(u\) is diagonal in the basis B and moreover the family \(\lambda\) of the eigenvalues of \(u\) belongs to \(\mathbf{R}_+^I\) and is summable. The family \(\lambda\) therefore belongs to \(\mathcal{C}_0(\mathrm{I};\mathrm{K})\) (TG, III, p. 38, prop. 1), and consequently the endomorphism \(u\) is compact (prop. 2 of IV, p. 148).

Let us prove \(b\) ). Let \((e_i)_{i \in \mathbf{J}}\) be an orthonormal basis of E extending the family B, so that \(u(e_i) = 0\) if \(i \in \mathbf{J} - \mathbf{I}\) . We obtain

\[
\operatorname{Tr} (u) = \sum_ {i \in \mathrm{J}} \langle e _ {i} \mid u (e _ {i}) \rangle = \sum_ {i \in \mathrm{I}} \langle e _ {i} \mid u (e _ {i}) \rangle = \sum_ {i \in \mathrm{I}} \alpha_ {i} \langle e _ {i} \mid f _ {i} \rangle .
\]

COROLLARY 2. --- Let F be a Hilbert space. Let \(u \in \mathcal{L}(\mathrm{E};\mathrm{F})\) be a Hilbert--Schmidt operator. Then u is a compact linear operator.

Let \((j,|u|)\) be the polar decomposition of u (I, p. 140, def. 4). By definition (EVT, V, p. 50, def. 9) the endomorphism \(u^{*}u\) has finite trace, hence is compact (corollary 1, a)); the same is true of \(|u|=\sqrt{u^{*}u}\) (prop. 6 of III, p. 91, b)) and of \(u=j|u|\) (prop. 3 of III, p. 5).

COROLLARY 3. --- Suppose that K = C. Let u be a positive compact endomorphism of E. The endomorphism u is of finite trace if and only if the decreasing family of its eigenvalues \((\lambda_{n}(u))_{n\in\mathrm{I}_{\mathrm{E}}}\) is summable; the trace of u is then the sum of this family.

By th. 1 of IV, p. 149, this follows from the preceding proposition and the definition of the sequence \((\lambda_n(u))_{n\in \mathrm{I_E}}\) (no. 3 of IV, p. 151), taking into account formula (28) of EVT, V, p. 49.

COROLLARY 4. --- Let F be a Hilbert space. Let \(u \in \mathcal{L}(\mathrm{E};\mathrm{F})\) be a compact linear map. Let \(\mathrm{B} = (e_i)_{i \in \mathrm{I}}\) be an orthonormal basis of the initial space \(\operatorname{Ker}(u)^{\circ}\) of \(u\) , \(\mathrm{C} = (f_i)_{i \in \mathrm{I}}\) be an orthonormal family of F and \((\alpha_{i})_{i\in\mathrm{I}}\) be a family in \((\mathbf{R}_{+}^{*})^{\mathrm{I}}\) such that \(u(e_{i})=\alpha_{i}f_{i}\) for all \(i\in I\) (cor. 2 of IV, p. 150).

The endomorphism \(|u|\) of E has finite trace if and only if the family \((\alpha_{i})_{i\in \mathrm{I}}\) is summable. We then have \(u\in \mathcal{L}_2(\mathrm{E};\mathrm{F})\) and

\[
\operatorname{Tr} (| u |) = \sum_ {i \in \mathrm{I}} \alpha_ {i}, \quad \| u \| _ {2} ^ {2} = \operatorname{Tr} (u ^ {*} u) = \sum_ {i \in \mathrm{I}} \alpha_ {i} ^ {2},\tag{8}
\]

and in particular \(\|u\|_{2}\leqslant\mathrm{Tr}(|u|)\) .

The family of nonzero eigenvalues of \(|u|\) is \((\alpha_{i})_{i\in \mathrm{I}}\) , hence \(|u|\) has finite trace if and only if the family \((\alpha_{i})_{i\in \mathrm{I}}\) is summable (prop. 12). If that is the case, the family \((\alpha_{i}^{2})\) of nonzero eigenvalues of \(u^{*}u = |u|^{2}\) is summable, and formulas (8) result from loc. cit.

Lemma 7. --- Let \(\lambda = (\lambda_i)_{i \in \mathrm{I}}\) be a family of complex numbers. For every \(t \in \mathbf{C}^*\) , let \(n_t\) be the cardinality of the set of \(i \in \mathrm{I}\) such that \(\lambda_i = t\) . The family \((\lambda_i)_{i \in \mathrm{I}}\) is summable if and only if \(n_t\) is finite for all \(t \in \mathbf{C}^*\) and if the family \((n_t t)_{t \in \mathbf{C}^*}\) is summable. In this case, the sums of these two families are equal.

Suppose that the family \((\lambda_{i})_{i\in\mathrm{I}}\) is summable. For every \(t\in\mathbf{C}\) , let \(I_{t}\) be the set of \(i\in I\) such that \(\lambda_{i}=t\) . By TG, III, p. 39, th. 2, applied to the partition of I by the sets \(I_{t}\) , the set \(I_{t}\) is finite for all \(t\in\mathbf{C}^{*}\) , and moreover the family \((n_{t}t)_{t\in\mathbf{C}^{*}}\) is summable with sum equal to that of the family \((\lambda_{i})_{i\in\mathrm{I}}\) .

Conversely, suppose that \(n_{t}\) is finite for all \(t \in \mathbf{C}^{*}\) and that the family \((n_{t}t)_{t \in \mathbf{C}^{*}}\) is summable. Let J be a finite subset of I and \(\Lambda\) the set of \(\lambda_{i}\) for \(i \in J\) . We have

\[
\left| \sum_ {i \in J} \lambda_ {i} \right| \leqslant \sum_ {t \in \Lambda - \{0 \}} n _ {t} | t | \leqslant \sum_ {t \in \mathbf {C} ^ {*}} n _ {t} | t |,
\]

therefore the family \((\lambda_{i})_{i\in\mathrm{I}}\) is summable (TG, VII, p. 17, corollary).

PROPOSITION 13. --- Suppose that K = C. Let u be a compact normal endomorphism of E. For \(t \in \mathrm{Sp}_{\mathrm{s}}(u)\) , let \(n_{t} \geqslant 1\) be the spectral multiplicity of the eigenvalue t of u. For u to have finite trace, it is necessary and sufficient that the family \((n_{t}t)_{t \in \mathrm{Sp}_{\mathrm{s}}(u)}\) be summable. The trace of u is then the sum of this family.

Let \(\mathrm{B} = (e_i)_{i \in \mathrm{I}}\) be an orthonormal basis of E such that u is diagonal in the basis B (th. 1 of IV, p. 149), and let \(\lambda = (\lambda_i)_{i \in \mathrm{I}}\) be the family of its eigenvalues. Since, for \(t \in \mathrm{Sp}_\mathrm{s}(u)\) , the multiplicity \(n_t\) is equal to the number of elements \(i \in I\) such that \(\lambda_i = t\), and since the nonzero elements of \(\lambda\) belong to \(\mathrm{Sp}_{\mathrm{s}}(u)\) , the assertion then follows from prop. 12 and the preceding lemma.

\subsection{Nuclear maps}

In this section, F denotes a Hilbert space over K. When K = C, we recall (I, p. 139, def. 3) that for \(u \in \mathcal{L}(\mathrm{E}; \mathrm{F})\) , we denote by \(|u|\) the positive endomorphism \(\sqrt{u^{*} \circ u}\) of E. When K = R, the element \(|u_{(\mathbf{C})}|\) of \(\mathcal{L}(\mathrm{E}_{(\mathbf{C})})\) is of the form \(v_{(\mathbf{C})}\) for a unique endomorphism \(v \in \mathcal{L}(\mathrm{E})\) , which is still denoted \(|u|\) (cf. I, p. 87).

We denote by \((u,v)\mapsto\langle u|v\rangle=\mathrm{Tr}(u^{*}v)\) the inner product in the Hilbert space \(\mathcal{L}_{2}(\mathrm{E};\mathrm{F})\) (EVT, V, p. 53, remarks 1 and 2).

For every \(u \in \mathcal{L}(\mathrm{E};\mathrm{F})\) , we denote by \(\| u\| _1\) the quantity \(\mathrm{Tr}(|u|)\) if \(|u|\) has finite trace and \(\| u\| _1 = +\infty\) if this is not the case. We therefore have \(\| u\| _1 \in \mathbf{R}_+ \cup \{+\infty\}\) . Since \(\| u\| = \| |u|\|\) (prop. 10, a) of I, p. 139) and since, if \(v\) is positive and of finite trace, then \(\mathrm{Tr}(v) \geqslant \| v\|\) (EVT, V, p. 49, formula (24bis) and p. 44, prop. 9), it follows that

\[
\| u \| \leqslant \| u \| _ {1}.\tag{9}
\]

If \(\| u\| _1\) is finite, then \(u\) is a Hilbert--Schmidt operator and \(\| u\| _2 \leqslant \| u\| _1\) (cor. 4 of IV, p. 165).

PROPOSITION 14. --- Let \(u \in \mathcal{L}_2(\mathrm{E};\mathrm{F})\) . We have

\[
\| u\|_{1} = \sup_{\substack{v\in \mathcal{L}_{2}(\mathrm{E};\mathrm{F})\\ \| v\| \leqslant 1}}|\langle v  |  u\rangle |,
\]

the supremum being computed in \(\mathbf{R}_{+} \cup \{+\infty\}\) .

Let \(\mathrm{B} = (e_i)_{i \in \mathrm{I}}\) be an orthonormal basis of the initial space \(\operatorname{Ker}(u)^{\circ}\) of u, \(\mathrm{C} = (f_i)_{i \in \mathrm{I}}\) be an orthonormal family of F and \((\alpha_i)_{i \in \mathrm{I}}\) be a family in \((\mathbf{R}_+^*)^{\mathrm{I}}\) such that \(u(e_i) = \alpha_i f_i\) for all \(i \in I\) (cor. 2 of IV, p. 150). The family \((\alpha_i)_{i \in \mathrm{I}}\) is square-integrable since u belongs to \(\mathcal{L}_2(\mathrm{E};\mathrm{F})\) (cor. 4 of IV, p. 165). Let \((e_j)_{j \in \mathrm{J}}\) be an orthonormal basis of E extending \((e_i)_{i \in \mathrm{I}}\) .

Let L be a finite subset of I. Let \(v_{L}\) be the continuous finite-rank linear map from E into F defined by

\[
v _ {\mathrm{L}} (x) = \sum_ {i \in \mathrm{L}} \langle e _ {i} | x \rangle f _ {i}
\]

for all \(x\) in E. We have \(\| v_{\mathrm{L}}\| \leqslant 1\) and \(v_{\mathrm{L}}\in \mathcal{L}_{2}(\mathrm{E};\mathrm{F})\) . Moreover, for every \(j\in \mathrm{J}\) , we have \(v_{\mathrm{L}}(e_j) = 0\) if \(j\notin \mathrm{L}\) and \(v_{\mathrm{L}}(e_j) = f_j\) if \(j\in \mathrm{L}\). Thus

\[
| \langle v _ {\mathrm{L}} | u \rangle | = | \operatorname{Tr} (v _ {\mathrm{L}} ^ {*} u) | = \left| \sum_ {j \in \mathrm{J}} \langle v _ {\mathrm{L}} (e _ {j}) | u (e _ {j}) \rangle \right| = \sum_ {j \in \mathrm{L}} \alpha_ {j}.
\]

From loc. cit., we deduce that

\[
\| u\|_{1} = \sup_{\text{L}}\sum_{j\in \text{L}}\alpha_{j}\leqslant \sup_{\substack{v\in \mathcal{L}_{2}(\text{E};\text{F})\\ \| v\| \leqslant 1}}|\langle v\mid u\rangle |\tag{10}
\]

in \(\mathbf{R}_{+}\cup\{+\infty\}\) .

This implies the equality of the proposition when \(\|u\|_{1}=+\infty\) .

Suppose that \(\|u\|_{1}\) is finite. For every \(i \in I\) , let \(p_{i}\) be the linear map from E to F such that \(p_{i}(e_{i}) = f_{i}\) and \(p_{i}(x) = 0\) for all \(x \in e_{i}^{\circ}\) . For all j and k in I, we have

\[
\langle p _ {j} \mid p _ {k} \rangle = \mathrm{Tr} (p _ {j} ^ {*} p _ {k}) = \sum_ {i \in \mathrm{J}} \langle p _ {j} (e _ {i}) \mid p _ {k} (e _ {i}) \rangle ,
\]

which is zero unless \(j = k\) , in which case this quantity equals \(\| f_k\|^2 = 1\) . The family \((p_i)_{i \in \mathrm{I}}\) is therefore orthonormal in \(\mathcal{L}_2(\mathrm{E};\mathrm{F})\) . Consequently, the family \((\alpha_i p_i)_{i \in \mathrm{I}}\) is summable in \(\mathcal{L}_2(\mathrm{E};\mathrm{F})\) ; its sum is equal to \(u\) since these two continuous linear maps coincide on the elements \(e_i\) for all \(i \in \mathrm{J}\) .

Let \(v \in \mathcal{L}_2(\mathrm{E};\mathrm{F})\) . For every \(i \in \mathrm{I}\) , it follows that

\[
\langle v \mid p _ {i} \rangle = \operatorname{Tr} (v ^ {*} p _ {i}) = \sum_ {j \in \mathrm{J}} \langle v (e _ {j}) \mid p _ {i} (e _ {j}) \rangle = \langle v (e _ {i}) \mid f _ {i} \rangle .
\]

If \(\|v\| \leqslant 1\) , we have

\[
\begin{array}{c} | \langle v | u \rangle | = \left| \langle v | \sum_ {i \in I} \alpha_ {i} p _ {i} \rangle \right| \leqslant \sum_ {i \in I} \alpha_ {i} | \langle v | p _ {i} \rangle | \\ = \sum_ {i \in I} \alpha_ {i} | \langle v (e _ {i}) | f _ {i} \rangle | \leqslant \sum_ {i \in I} \alpha_ {i} = \mathrm{Tr} (| u |), \end{array}
\]

whence the inequality

\[
\sup_{\substack{v\in \mathcal{L}_{2}(\mathrm{E};\mathrm{F})\\ \| v\| \leqslant 1}}|\langle v\mid u\rangle |\leqslant \| u\|_{1},
\]

which, combined with formula (10), concludes the proof of the proposition.

It follows from this proposition that the set \(\mathcal{L}_1(\mathrm{E};\mathrm{F})\) of continuous linear maps \(u\) from E to F such that \(\| u\| _1\) is finite is a vector subspace of \(\mathcal{L}_{2}(\mathrm{E};\mathrm{F})\) and that the map \(u\mapsto \| u\| _1\) is a semi-norm on \(\mathcal{L}_1(\mathrm{E};\mathrm{F})\) ; inequality (9) shows that it is a norm.

DEFINITION 4. --- We say that the vector space \(\mathcal{L}_{1}(\mathrm{E};\mathrm{F})\) equipped with the norm \(u\mapsto \| u\|_{1}\) is the space of nuclear maps from E to F. If \(u\in \mathcal{L}_1(\mathrm{E};\mathrm{F})\) , one says that \(u\) is nuclear.

Remarks. --- 1) Suppose that K = R. Let \(u \in \mathcal{L}(E; F)\) . We have \(u \in \mathcal{L}_{1}(E; F)\) if and only if \(u_{(\mathbf{C})} \in \mathcal{L}_{1}(E_{(\mathbf{C})}; F_{(\mathbf{C})})\) ; in this case, we have \(\|u\|_{1} = \|u_{(\mathbf{C})}\|_{1}\) .

\begin{enumerate}
\def\labelenumi{\arabic{enumi})}
\setcounter{enumi}{1}
\tightlist
\item
  Let \(u\) be a nuclear map from E to F. The map \(u\) being Hilbert--Schmidt, it is compact (cor. 2 of IV, p. 165). Moreover, proposition 14 implies that for every \(v \in \mathcal{L}_2(\mathrm{E};\mathrm{F})\) , we have
\end{enumerate}

\[
| \langle v | u \rangle | \leqslant \| v \| \| u \| _ {1}.\tag{11}
\]

\begin{enumerate}
\def\labelenumi{\arabic{enumi})}
\setcounter{enumi}{2}
\item
  Suppose that \(\mathrm{K} = \mathbf{C}\). Let \(u \in \mathcal{L}_1(\mathrm{E};\mathrm{E})\) such that \(u\) is positive. The norm \(\| u\| _1\) of \(u\) is the sum of the sequence \((\lambda_n(u))_{n\in \mathrm{I_E}}\) (cor. 3 of IV, p. 165).
\item
  The canonical inclusion of \(\mathcal{L}_1(\mathrm{E};\mathrm{F})\) in \(\mathcal{L}(\mathrm{E};\mathrm{F})\) has norm \(\leqslant 1\) (inequality (9), p. 167).
\end{enumerate}

PROPOSITION 15. --- Let G be a Hilbert space over K. The map \((u,v)\mapsto v\circ u\) from \(\mathcal{L}(\mathrm{E};\mathrm{F})\times\mathcal{L}(\mathrm{F};\mathrm{G})\) into \(\mathcal{L}(\mathrm{E};\mathrm{G})\) defines by passage to quotients a continuous bilinear map of norm \(\leqslant1\) from \(\mathcal{L}_{2}(\mathrm{E};\mathrm{F})\times\mathcal{L}_{2}(\mathrm{F};\mathrm{G})\) into \(\mathcal{L}_{1}(\mathrm{E};\mathrm{G})\) .

Let \(u \in \mathcal{L}_2(\mathrm{E};\mathrm{F})\) and \(v \in \mathcal{L}_2(\mathrm{F};\mathrm{G})\). Let \(w \in \mathcal{L}_2(\mathrm{E};\mathrm{G})\) . We obtain \(\langle uv | w \rangle = \operatorname{Tr}(v^* u^* w) = \langle v | u^* w \rangle\) whence

\[
| \langle u v | w \rangle | \leqslant \| v \| _ {2} \| u ^ {*} w \| _ {2} \leqslant \| v \| _ {2} \| u \| _ {2} \| w \|
\]

by the Cauchy--Schwarz inequality and formula (37) of EVT, V, p. 52. The result therefore follows from prop. 14.

Lemma 8. --- The map \(u \mapsto u^{*}\) from \(\mathcal{L}(\mathrm{E};\mathrm{F})\) into \(\mathcal{L}(\mathrm{F};\mathrm{E})\) defines by passage to quotients an isometric linear map from \(\mathcal{L}_{1}(\mathrm{E};\mathrm{F})\) into \(\mathcal{L}_{1}(\mathrm{F};\mathrm{E})\) .

Let \(u \in \mathcal{L}_1(\mathrm{E};\mathrm{F})\) . Since \(u\) is a Hilbert--Schmidt map, so is \(u^*\) (EVT, V, p. 54). The map \(v \mapsto v^*\) is a bijection from the set of \(v \in \mathcal{L}_2(\mathrm{E};\mathrm{F})\) such that \(\|v\| \leqslant 1\) onto the set of \(w \in \mathcal{L}_2(\mathrm{F};\mathrm{E})\) such that \(\|w\| \leqslant 1\) ; for every \(v \in \mathcal{L}_2(\mathrm{E};\mathrm{F})\)

with \(\|v\|\leqslant1\) , we have \(\langle v|u\rangle=\langle u^{*}|v^{*}\rangle\) (EVT, V, p. 54, formula (42)), whence the result by prop. 14.

PROPOSITION 16. --- Let \(E_{1}\) and \(F_{1}\) be Hilbert spaces. Let \(u\) be in \(\mathcal{L}_{1}(\mathrm{E};\mathrm{F})\), \(v\) be in \(\mathcal{L}(\mathrm{E}_{1};\mathrm{E})\) and \(v_{2}\) be in \(\mathcal{L}(\mathrm{F};\mathrm{F}_{1})\). We then have \(v_{2}uv_{1} \in \mathcal{L}_{1}(\mathrm{E}_{1};\mathrm{F}_{1})\) and \(\|v_{2}uv_{1}\|_{1} \leqslant \|v_{2}\|\|v_{1}\|\|u\|_{1}\) .

Let \(w \in \mathcal{L}_2(\mathrm{E};\mathrm{F}_1)\) such that \(\| w\| \leqslant 1\) . We have \(v_{2}u \in \mathcal{L}_{2}(\mathrm{E};\mathrm{F}_{1})\) (EVT, V, p. 52, formula (36)). Since \(v_{2}^{*}w \in \mathcal{L}_{2}(\mathrm{E};\mathrm{F})\) (loc. cit.), it follows

\[
| \langle w | v _ {2} u \rangle | = | \langle v _ {2} ^ {*} w | u \rangle | \leqslant \| v _ {2} \| \| u \| _ {1}
\]

(formula (11)), whence \(v_2u \in \mathcal{L}_1(\mathrm{E};\mathrm{F}_1)\) and \(\| v_2u\| _1 \leqslant \| v_2\| \| u\| _1\) (prop. 14). Let \(v_{1} \in \mathcal{L}(\mathrm{E}_{1};\mathrm{E})\) ; since \(uv_{1} = (v_{1}^{*}u^{*})^{*}\) , we have \(uv_{1} \in \mathcal{L}_{1}(\mathrm{E}_{1};\mathrm{F})\) and \(\| uv_{1}\|_{1} \leqslant \| v_{1}^{*}\| \| u^{*}\|_{1} = \| v_{1}\| \| u\|_{1}\) (lemma 8).

The proposition follows immediately from these inequalities.

PROPOSITION 17. — The space \(\mathcal{L}_{1}(\mathrm{E};\mathrm{E})\) coincides with the space \(\mathcal{L}_{1}(\mathrm{E})\) of endomorphisms of finite trace of E, and we have \(|\operatorname{Tr}(u)|\leqslant \| u\|_{1}\) for every endomorphism \(u\) of E with finite trace.

We may assume that K = C. The space \(\mathcal{L}_{1}(E;E)\) contains by definition the set of positive endomorphisms of finite trace, and hence \(\mathcal{L}_{1}(E)\subset\mathcal{L}_{1}(E;E)\) (EVT, p. 50, def. 8). Conversely, let us show that \(\mathcal{L}_{1}(E;E)\) is contained in \(\mathcal{L}_{1}(E)\) .

Let \(u \in \mathcal{L}_{1}(\mathrm{E};\mathrm{E})\) . We have \(u^{*} \in \mathcal{L}_{1}(\mathrm{E};\mathrm{E})\) (lemma 8); it therefore suffices to show that the Hermitian elements of \(\mathcal{L}_{1}(\mathrm{E};\mathrm{E})\) are of finite trace (lemma 2 of I, p. 96).

Let u be such an endomorphism. It is compact (remark 2, p. 169). Let B be an orthonormal basis of E such that u is diagonal in the basis B (th. 1 of IV, p. 149), and let \(\lambda\) be the family of eigenvalues of u relative to B. The endomorphism \(|u|\) is diagonalizable in the basis B and the family of its eigenvalues is \(|\lambda|\) (prop. 1, d) of IV, p. 147). Since \(|u|\) is of finite trace, this latter family is summable (prop. 12 of IV, p. 164), hence the family \(\lambda\) is summable. Consequently, u is of finite trace (loc. cit.), and we have \(|\mathrm{Tr}(u)| \leqslant \mathrm{Tr}(|u|) = \|u\|_{1}\) .

Consequently, the space \(\mathcal{L}_{1}(\mathrm{E})\) is a self-adjoint two-sided ideal of the C*-algebra \(\mathcal{L}(\mathrm{E})\) . If E is of infinite dimension, this ideal is not closed in \(\mathcal{L}(\mathrm{E})\) .

Propositions 17 and 16 show that the map \(b\) from \(\mathcal{L}_1(\mathrm{E};\mathrm{F})\times \mathcal{L}(\mathrm{F};\mathrm{E})\) into \(\mathbf{C}\) defined by \(b(u,v) = \operatorname {Tr}(vu)\) is a continuous bilinear form which puts the spaces \(\mathcal{L}_1(\mathrm{E};\mathrm{F})\) and \(\mathcal{L}(\mathrm{F};\mathrm{E})\) in duality.

Lemma 9. --- Let \(u \in \mathcal{L}_1(\mathrm{E};\mathrm{F})\) . We have

\[
\| u\|_{1} = \sup_{\substack{v\in \mathcal{L}^{c}(\mathrm{F};\mathrm{E})\\ \| v\| \leqslant 1}}|b(u,v)|.
\]

Since every Hilbert--Schmidt map is compact (cor. 2 of IV, p. 165), we have

\[
\| u\|_{1}\leqslant \sup_{\substack{v\in \mathcal{L}^{\mathrm{c}}(\mathrm{E};\mathrm{F})\\ \| v\| \leqslant 1}}|\operatorname{Tr}(vu)| = \sup_{\substack{v\in \mathcal{L}^{\mathrm{c}}(\mathrm{F};\mathrm{E})\\ \| v\| \leqslant 1}}|b(u,v)|
\]

by prop. 14.

Let \(v \in \mathcal{L}^{c}(\mathrm{E};\mathrm{F})\) be of norm \(\leqslant 1\) . Since \(\mathcal{L}_{2}(\mathrm{E};\mathrm{F})\) is dense in \(\mathcal{L}^{c}(\mathrm{E};\mathrm{F})\) , there exists a sequence \((v_{n})_{n\in\mathbf{N}}\) in \(\mathcal{L}_{2}(\mathrm{E};\mathrm{F})\) which converges to v in \(\mathcal{L}(\mathrm{E};\mathrm{F})\) . The sequence \((b(u,v_{n}))_{n\in\mathbf{N}}\) converges to \(b(u,v)\) ; since \(|b(u,v_{n})| = |\operatorname{Tr}(v_{n}u)| = |\langle v_{n}^{*}|u\rangle| \leqslant \|u\|_{1}\) (prop. 14) for all \(n \in N\) , it follows that \(|b(u,v)| \leqslant \|u\|_{1}\) .

Denote by \(\theta\) the continuous linear map

\[
\theta \colon \mathcal {L} _ {1} (\mathrm{E}; \mathrm{F}) \to \mathcal {L} ^ {\mathrm{c}} (\mathrm{F}; \mathrm{E}) ^ {\prime}
\]

such that \(\theta(u)(v) = b(u,v)\) .

PROPOSITION 18. --- The map θ is an isometric isomorphism.

By lemma 9, the map \(\theta\) is isometric, and it suffices to show that \(\theta\) is surjective.

Let \(\lambda \in \mathcal{L}^{\mathrm{c}}(\mathrm{F};\mathrm{E})'\) . Since \(\| v\| \leqslant \| v\|_{2}\) for all \(v\in \mathcal{L}_2(\mathrm{F};\mathrm{E})\) (EVT, V, p. 52, formula (33)), the restriction of \(\lambda\) to the subspace \(\mathcal{L}_2(\mathrm{F};\mathrm{E})\) is a continuous linear form on the Hilbert space \(\mathcal{L}_2(\mathrm{F};\mathrm{E})\) . There therefore exists an element \(u\) of \(\mathcal{L}_2(\mathrm{F};\mathrm{E})\) such that \(\lambda (v) = \langle u|v\rangle\) for all \(v\in \mathcal{L}_2(\mathrm{F};\mathrm{E})\) (EVT, V, p. 15, th. 3).

For every \(w \in \mathcal{L}_2(\mathrm{E};\mathrm{F})\) of norm \(\leqslant 1\) , we have \(|\langle w|u\rangle| = |\lambda(w)| \leqslant \| \lambda \|\) . Consequently, \(u\) is a nuclear map from \(\mathrm{E}\) into \(\mathrm{F}\) (prop. 14).

The continuous linear forms \(\lambda\) and \(\theta(u)\) are equal on the dense subspace \(\mathcal{L}_{2}(\mathrm{F},\mathrm{E})\) of \(\mathcal{L}^{\mathrm{c}}(\mathrm{F};\mathrm{E})\) . It follows that \(\lambda=\theta(u)\) , which concludes the proof.

COROLLARY. --- The normed space \(\mathcal{L}_{1}(\mathrm{E};\mathrm{F})\) is a Banach space.

The assertion follows from the proposition and from EVT, III, p. 24, cor. 2, since the space \(\mathcal{L}^{\mathrm{c}}(\mathrm{F};\mathrm{E})\) is a normed space.

\subsection{Hilbert--Schmidt integral operators}

In this issue, X and Y are locally compact topological spaces. Let \(\mu\) be a positive measure on X and \(\nu\) be a positive measure on Y. We denote by \(\mathrm{L}^{2}(\mathrm{X})\) , \(\mathrm{L}^{2}(\mathrm{Y})\) and \(\mathrm{L}^{2}(\mathrm{X} \times \mathrm{Y})\) the spaces \(\mathrm{L}^{2}(\mathrm{X}, \mu)\) , \(\mathrm{L}^{2}(\mathrm{Y}, \nu)\) and \(\mathrm{L}^{2}(\mathrm{X} \times \mathrm{Y}, \mu \otimes \nu)\) respectively, and likewise for \(\mathcal{L}^{2}(\mathrm{X})\) , \(\mathcal{L}^{2}(\mathrm{Y})\) and \(\mathcal{L}^{2}(\mathrm{X} \times \mathrm{Y})\) .

We recall (cf. No. 4 of III, p. 26) that for every \(N \in \mathcal{L}^{2}(X \times Y)\) there exists a unique continuous linear map \(u_{\mathrm{N}}\) from \(L^{2}(X)\) into \(L^{2}(Y)\) such that

\[
\langle g \mid u _ {\mathrm{N}} (f) \rangle = \int_ {\mathrm{X} \times \mathrm{Y}} \overline {{g (y)}} \mathrm{N} (x, y) f (x) d (\mu \otimes \nu) (x, y)\tag{12}
\]

for all \(f \in \mathrm{L}^{2}(\mathrm{X})\) and every \(g \in \mathrm{L}^{2}(\mathrm{Y})\) . The map \(u_{N}\) is compact (cor. 1 of III, p. 33). The map \(N \mapsto u_{N}\) induces, by passage to quotients, a continuous and injective linear map from \(\mathrm{L}^{2}(\mathrm{X} \times \mathrm{Y})\) into \(\mathcal{L}(\mathrm{L}^{2}(\mathrm{X});\mathrm{L}^{2}(\mathrm{Y}))\) (prop. 5 of III, p. 30).

We denote by \(\theta\) the unique linear map from \(\mathcal{L}^2 (\mathrm{X})\otimes \mathcal{L}^2 (\mathrm{Y})\) into \(\mathcal{L}^2 (\mathrm{X}\times \mathrm{Y})\) which, for all \(f\in \mathcal{L}^2 (\mathrm{X})\) and \(g\in \mathcal{L}^2 (\mathrm{Y})\), sends \(f\otimes g\) to the function defined by \((x,y)\mapsto f(x)g(y)\) .

Lemma 10. --- The map \(\theta\) defines by passage to quotients a linear map \(\widetilde{\theta}\) from \(\mathrm{L}^2 (\mathrm{X})\otimes \mathrm{L}^2 (\mathrm{Y})\) into \(\mathrm{L}^2 (\mathrm{X}\times \mathrm{Y})\) and there exists a unique isometric isomorphism from \(\mathrm{L}^2 (\mathrm{X})\widehat{\otimes}_2\mathrm{L}^2 (\mathrm{Y})\) onto \(\mathrm{L}^2 (\mathrm{X}\times \mathrm{Y})\) which coincides with it on \(\mathrm{L}^2 (\mathrm{X})\otimes \mathrm{L}^2 (\mathrm{Y})\) .

The first assertion is elementary. Let us prove the second.

Let \((f_i)_{i\in \mathrm{I}}\) and \((g_j)_{j\in \mathrm{J}}\) be orthonormal bases of \(\mathrm{L}^2 (\mathrm{X})\) and \(\mathrm{L}^2 (\mathrm{Y})\) respectively. The family \((\overline{f}_i\otimes g_j)_{(i,j)\in \mathrm{I}\times \mathrm{J}}\) is an orthonormal basis of \(\mathrm{L}^2 (\mathrm{X})\widehat{\otimes}_2\mathrm{L}^2 (\mathrm{Y})\) (EVT, V, p. 29, cor. 1) and the family \((\widetilde{\theta} (\overline{f}_i\otimes g_j))_{(i,j)\in \mathrm{I}\times \mathrm{J}}\) is orthonormal in \(\mathrm{L}^2 (\mathrm{X}\times \mathrm{Y})\). The map \(\widetilde{\theta}\) extends therefore by continuity to an isometric linear map from \(\mathrm{L}^2 (\mathrm{X})\widehat{\otimes}_2\mathrm{L}^2 (\mathrm{Y})\) into \(\mathrm{L}^2 (\mathrm{X}\times \mathrm{Y})\). It remains to show that this extension is surjective.

For this, it suffices to prove that the image of \(\widetilde{\theta}\) is dense in \(\mathrm{L}^{2}(\mathrm{X}\times\mathrm{Y})\). Let N be an element of \(\mathrm{L}^{2}(\mathrm{X}\times\mathrm{Y})\) orthogonal to the image of \(\widetilde{\theta}\) . For all \(i\in I\) and \(j\in J\) , we have \(\langle g_{j}\mid u_{\mathrm{N}}(f_{i})\rangle=\langle\widetilde{\theta}(\overline{f}_{i}\otimes g_{j})\mid\mathrm{N}\rangle=0\). From formula (12), it follows that \(u_{N}=0\) and therefore N=0.

The preceding lemma establishes the assertion stated in EVT, V, p. 29, example 2.

In what follows, we will identify \(\mathrm{L}^{2}(\mathrm{X})\widehat{\otimes}_{2}\mathrm{L}^{2}(\mathrm{Y})\) and \(\mathrm{L}^{2}(\mathrm{X}\times\mathrm{Y})\) by the isometric isomorphism of Lemma 10.

PROPOSITION 19. --- The map \(N\mapsto u_{\mathrm{N}}\) is an isometric isomorphism from \(\mathrm{L}^{2}(\mathrm{X}\times\mathrm{Y})\) onto the space \(\mathcal{L}_{2}(\mathrm{L}^{2}(\mathrm{X});\mathrm{L}^{2}(\mathrm{Y}))\) of Hilbert--Schmidt maps from \(\mathrm{L}^{2}(\mathrm{X})\) into \(\mathrm{L}^{2}(\mathrm{Y})\) .

The linear map from \(\mathrm{L}^{2}(\mathrm{Y})\otimes\overline{\mathrm{L}^{2}(\mathrm{X})}\) into \(\mathcal{L}_{2}(\mathrm{L}^{2}(\mathrm{X});\mathrm{L}^{2}(\mathrm{Y}))\) that sends \(g\otimes f\) to the Hilbert--Schmidt map \(h\mapsto\langle f|h\rangle g\) extends to an isometric isomorphism \(\theta_{1}\) of \(\mathrm{L}^{2}(\mathrm{Y})\widehat{\otimes}_{2}\overline{\mathrm{L}^{2}(\mathrm{X})}\) onto \(\mathcal{L}_{2}(\mathrm{L}^{2}(\mathrm{X});\mathrm{L}^{2}(\mathrm{Y}))\) (EVT, V, p. 52, th. 1).

Additionally, let us denote by \(\theta_{2}\) the isometric isomorphism of \(\mathrm{L}^{2}(\mathrm{X})\widehat{\otimes}_{2}\mathrm{L}^{2}(\mathrm{Y})\) onto \(\mathrm{L}^{2}(\mathrm{Y})\widehat{\otimes}_{2}\overline{\mathrm{L}^{2}(\mathrm{X})}\) which sends \(f\otimes g\) to \(g\otimes \overline{f}\) for all \(f\in \mathrm{L}^2 (\mathrm{X})\) and every \(g\in \mathrm{L}^2 (\mathrm{Y})\) .

The linear map \(\theta_{3} = \theta_{1} \circ \theta_{2}\) is identified with an isometric isomorphism of \(\mathrm{L}^{2}(\mathrm{X} \times \mathrm{Y})\) onto \(\mathcal{L}_{2}(\mathrm{L}^{2}(\mathrm{X}); \mathrm{L}^{2}(\mathrm{Y}))\) .

Let \(f \in \mathrm{L}^2(\mathrm{X})\) and \(g \in \mathrm{L}^2(\mathrm{Y})\) and let \(\mathrm{N}\) be the element of \(\mathrm{L}^2(\mathrm{X} \times \mathrm{Y})\) identified with \(f \otimes g\). By INT, V, p. 95, § 8, no. 3, cor. 2 and formula (12), we have

\[
\langle g _ {1} \mid u _ {\mathrm{N}} (f _ {1}) \rangle = \langle \overline {{f}} \mid f _ {1} \rangle \langle g _ {1} \mid g \rangle
\]

for all \(f_1 \in \mathrm{L}^2(\mathrm{X})\) and \(g_1 \in \mathrm{L}^2(\mathrm{Y})\) . The map \(u = \theta_3(\mathrm{N})\) is the linear map \(\theta_1(g \otimes \overline{f})\); it therefore satisfies \(u(h) = \langle \overline{f} | h \rangle g\) for all \(h \in \mathrm{L}^2(\mathrm{X})\). Consequently, for all \(f_1 \in \mathrm{L}^2(\mathrm{X})\) and \(g_1 \in \mathrm{L}^2(\mathrm{Y})\) , it follows that

\[
\langle g _ {1} \mid u (h _ {1}) \rangle = \langle \overline {{f}} \mid h _ {1} \rangle \langle g _ {1} \mid g \rangle = \langle g _ {1} \mid u _ {\mathrm{N}} (h _ {1}) \rangle ,
\]

whence \(\theta_3(N) = u_N\) . Since \(\theta_3\) and \(N \mapsto u_N\) are continuous, we conclude that \(\theta_3(N) = u_N\) for all \(N \in L^2(X \times Y)\) , which concludes the proof.

COROLLARY. --- For every \(N \in L^{2}(X \times Y)\) , the linear map \(u_{N}\) is a Hilbert--Schmidt operator and we have \(\mathrm{Tr}(u_{\mathrm{N}}^{*}u_{\mathrm{N}}) = \|N\|^{2}\) .

Indeed, one has \(\| u\| _2 = \sqrt{\mathrm{Tr}(u^*u)}\) for all \(u\in \mathcal{L}_2(\mathrm{L}^2 (\mathrm{X});\mathrm{L}^2 (\mathrm{Y}))\) .

Note. --- Let \(N \in L^{2}(X \times Y)\) . By Corollary 2 of IV, p. 150, there exist a countable set I, orthonormal families \((f_{i})_{i \in I}\) in \(L^{2}(X)\) and \((g_{i})_{i \in I}\) in \(L^{2}(Y)\) , and a family \((\alpha_{i})_{i \in I}\) in \(\mathbf{R}_{+}^{*}\) such that

\[
u _ {\mathrm{N}} (f) = \sum_ {i \in \mathrm{I}} \alpha_ {i} \langle f _ {i} | f \rangle g _ {i}
\]

for all \(f \in \mathrm{L}^2(\mathrm{X})\), where the series converges in \(\mathrm{L}^2(\mathrm{Y})\) . By cor. 4 of IV, p. 165 and the corollary above, we have

\[
\sum_ {i \in \mathrm{I}} \alpha_ {i} ^ {2} = \mathrm{Tr} (u _ {\mathrm{N}} ^ {*} u _ {\mathrm{N}}) = \| u _ {\mathrm{N}} \| _ {2} ^ {2} = \int_ {\mathrm{X} \times \mathrm{Y}} | \mathrm{N} (x, y) | ^ {2} d (\mu \otimes \nu) (x, y).
\]

Additionally, let us note that \(h_{i,j} = \overline{f}_{i} \otimes g_{j} \in \mathrm{L}^{2}(\mathrm{X} \times \mathrm{Y})\) for all \((i, j) \in \mathrm{I} \times \mathrm{I}\) . We then have

\[
\mathrm{N} = \sum_ {i \in \mathrm{I}} \alpha_ {i} h _ {i, i}
\]

in \(\mathrm{L}^2 (\mathrm{X}\times \mathrm{Y})\) .

Indeed, let \((f_{j})_{j\in\mathrm{J}}\) and \((g_{k})_{k\in\mathrm{K}}\) be orthonormal bases of \(\mathrm{L}^{2}(\mathrm{X})\) and \(\mathrm{L}^{2}(\mathrm{Y})\) , respectively, extending the families \((f_{i})_{i\in\mathrm{I}}\) and \((g_{i})_{i\in\mathrm{I}}\) . Set \(h_{j,k}=\overline{f}_{j}\otimes g_{k}\) for all \((j,k)\in\mathrm{J}\times\mathrm{K}\) . By Lemma 10, the family \((h_{j,k})_{(j,k)\in\mathrm{J}\times\mathrm{K}}\) is an orthonormal basis of \(\mathrm{L}^{2}(\mathrm{X}\times\mathrm{Y})\) . We have \(\langle h_{j,k}\mid\mathrm{N}\rangle=\langle g_{k}\mid u_{\mathrm{N}}(f_{j})\rangle\) for all \((j,k)\in\mathrm{J}\times\mathrm{K}\) . If \(j\notin I\) this quantity is zero. If \(j\in I\) , it is equal to \(\alpha_{j}\langle g_{k}\mid g_{j}\rangle\) , hence is zero unless k=j, in which case \(\langle h_{j,j}\mid N\rangle=\alpha_{j}\) . Consequently

\[
\mathrm{N} = \sum_ {(j, k) \in \mathrm{J} \times \mathrm{K}} \left\langle h _ {j, k} \mid \mathrm{N} \right\rangle h _ {j, k} = \sum_ {i \in \mathrm{I}} \alpha_ {i} h _ {i, i}.
\]

\subsection{Trace of integral operators with continuous kernel}

In this section, we keep the conventions of the preceding section with Y = X and \(\nu = \mu\). We assume that X is a locally compact topological space that is countable at infinity (TG, I, p. 68, def. 5).

In particular, we identify the spaces \(\mathrm{L}^{2}(\mathrm{X}\times\mathrm{X})\) and \(\mathrm{L}^{2}(\mathrm{X})\widehat{\otimes}_{2}\mathrm{L}^{2}(\mathrm{X})\) (lemma 10 of IV, p. 172). We shall denote by \(\widetilde{f}\) the class in \(\mathrm{L}^{2}(\mathrm{X})\) (resp. in \(\mathrm{L}^{2}(\mathrm{X}\times\mathrm{X})\) ) of a function \(f\in\mathcal{L}^{2}(\mathrm{X})\) (resp. in \(\mathcal{L}^{2}(\mathrm{X}\times\mathrm{X})\) ).

PROPOSITION 20. --- Let \((f_n)_{n \in \mathbf{N}}\) be a sequence of measurable maps from X into a metrizable space F. Suppose that the limit of \(f_n(x)\) exists outside a \(\mu\)-negligible subset of X. There exists a sequence \((\mathrm{C}_m)_{m \in \mathbf{N}}\) of compact subsets of X whose union \(\widetilde{\mathrm{X}}\) satisfies \(|\mu|(\mathrm{X} - \widetilde{\mathrm{X}}) = 0\) , such that the functions \(f_n\) are continuous on \(\mathrm{C}_m\) for all \(n \in \mathbf{N}\) and every \(m \in \mathbf{N}\) , and the sequence \((f_n)_{n \in \mathbf{N}}\) converges uniformly to \(f\) on \(\mathrm{C}_m\) for all \(m \in \mathbf{N}\) .

It follows from th. 2 of INT, IV, p. 175, § 5, no. 4 and from prop. 12, b) of INT, IV, p. 188, § 5, no. 8, that the set of compact subsets C of X such that the functions \(f_n\) are continuous on C for every \(n\) , and such that \((f_n)\) converges uniformly to \(f\) on C, is \(\mu\) -dense in X (INT, IV, p. 189, § 5, no. 8, def. 6). The assertion then follows from the remark of INT, IV, p. 189, § 5, no. 8.

Let \(N\) be a function belonging to \(\mathcal{L}^2 (\mathrm{X}\times \mathrm{X})\) and \(u_{\mathrm{N}}\) the Hilbert--Schmidt map from \(\mathrm{L}^2 (\mathrm{X})\) into \(\mathrm{L}^2 (\mathrm{X})\) with kernel \(N\) (Prop. 19 of IV, p. 173). Since the map \(u_{\mathrm{N}}\) is compact, there exist, by Prop. 9 of IV, p. 156, an element \(M\) of \(\overline{\mathbf{N}}\) and, denoting by \(I\) the set of integers \(n\in \mathbf{N}\) such that \(n\leqslant \mathrm{M}\), orthonormal families \((f_i)_{i\in \mathrm{I}}\) and \((g_{i})_{i\in \mathrm{I}}\) in \(\mathcal{L}^2 (\mathrm{X})\) and a family \((\alpha_{i})_{i\in \mathrm{I}}\) in \(\mathbf{R}_+^*\) such that

\[
u _ {\mathrm{N}} (f) = \sum_ {i \in \mathrm{I}} \alpha_ {i} \langle \widetilde {f} _ {i} | f \rangle \widetilde {g} _ {i}\tag{13}
\]

for all \(f \in \mathrm{L}^2(\mathrm{X})\), where the series converges in \(\mathrm{L}^2(\mathrm{X})\) .

For every \(i \in I\) , we denote by \(h_i \in \mathcal{L}^2(\mathrm{X} \times \mathrm{X})\) the function defined by \(h_i(x, y) = \overline{f_i(x)} g_i(y)\) , so that \(\widetilde{h}_i\) is the class of \(\overline{f}_i \otimes g_i\). By Remark 9 of IV, p. 173, the series with general term \(\alpha_i h_i\) converges to N in \(\mathrm{L}^2(\mathrm{X} \times \mathrm{X})\) .

In the rest of this issue, we assume that N is a continuous function.

PROPOSITION 21. --- Suppose that \(u_{\mathrm{N}}\) has finite trace. There then exists a set \(\widetilde{\mathrm{X}} \subset \mathrm{X}\) whose complement \(\mathrm{X} - \widetilde{\mathrm{X}}\) is \(\mu\) -negligible and a function \(\mathrm{H} \in \mathcal{L}^2(\mathrm{X} \times \mathrm{X})\) satisfying the following conditions:

\begin{enumerate}
\def\labelenumi{(\roman{enumi})}
\item
  For every \((x,y)\in\widetilde{\mathrm{X}}\times\widetilde{\mathrm{X}}\) the family \((\alpha_{i}\overline{f_{i}(x)}g_{i}(y))_{i\in\mathrm{I}}\) is summable in \(\mathbf{C}\) and its sum is \(\mathrm{N}(x,y)\) ;
\item
  For every finite subset J of I and every \((x,y)\in\widetilde{\mathbf{X}}\times\widetilde{\mathbf{X}}\) , we have
\end{enumerate}

\[
\left| \sum_ {i \in \mathrm{J}} \alpha_ {i} h _ {i} (x, y) \right| \leqslant \mathrm{H} (x, y);\tag{14}
\]

\begin{enumerate}
\def\labelenumi{(\roman{enumi})}
\setcounter{enumi}{2}
\tightlist
\item
  The function \(x \mapsto \mathrm{H}(x, x)\) belongs to \(\mathcal{L}^1(\mathrm{X})\) .
\end{enumerate}

Since \(u_{N}\) has finite trace, the family \((\alpha_{i})\) is summable (cor. 4 of IV, p. 165). The series

\[
\sum_ {i \in \mathrm{I}} \alpha_ {i} | f _ {i} | ^ {2}, \quad \sum_ {i \in \mathrm{I}} \alpha_ {i} | g _ {i} | ^ {2}
\]

therefore converge in \(\mathcal{L}^{1}(X)\) and consequently \(\mu\)-almost everywhere. Let F and G denote, respectively, the functions defined \(\mu\)-almost everywhere by the sum of these series, setting \(F(x) = 0\) and \(G(x) = 0\) for all \(x\) such that the corresponding series does not converge. Since F and G belong to \(\mathcal{L}^{1}(X)\), the function H defined by \(H(x,y) = \sqrt{F(x)G(y)}\) belongs to \(\mathcal{L}^{2}(X \times X)\) (INT, V, p. 95, § 8, no. 3, cor. 2).

Let us apply Prop. 20 to the space X, with \(F = \mathbf{C}^{2}\) and with the measurable maps

\[
s_{n}\colon x\mapsto \Bigl (\sum_{\substack{i\in \mathrm{I}\\ i\leqslant n}}\alpha_{i}|f_{i}(x)|^{2},\sum_{\substack{i\in \mathrm{I}\\ i\leqslant n}}\alpha_{i}|g_{i}(x)|^{2}\Bigr)
\]

defined on X. There therefore exists a sequence \((\mathrm{C}_{m})_{m\in\mathbf{N}}\) of compact subsets of X satisfying the following conditions:

\begin{enumerate}
\def\labelenumi{(\arabic{enumi})}
\item
  the union \(\widetilde{X}\) satisfies \(\mu(\mathrm{X}-\widetilde{\mathrm{X}})=0\) ;
\item
  for all \(m \in \mathbf{N}\) and every \(n \in \mathbf{N}\) , the functions \(s_n\) are continuous on \(\mathrm{C}_m\) ;
\item
  for all \(m \in \mathbf{N}\) , the series \(\sum \alpha_{i}|f_{i}|^{2}\) and \(\sum \alpha_{i}|g_{i}|^{2}\) converge uniformly on \(C_{m}\) to F and G, respectively;
\item
  the support of the measure induced by \(\mu\) on \(C_{m}\) is equal to \(C_{m}\) (possibly after replacing \(C_{m}\) by the support of this measure).
\end{enumerate}

Let us prove that the set \(\widetilde{\mathbf{X}}\) and the function H satisfy conditions (i), (ii), and (iii).

Let \((x,y)\in \widetilde{\mathbf{X}}\times \widetilde{\mathbf{X}}\). For every finite subset \(\mathrm{J}\subset \mathrm{I}\) , we have

\[
\left| \sum_ {i \in \mathrm{J}} \alpha_ {i} \overline {{f _ {i} (x)}} g _ {i} (y) \right| \leqslant \left(\sum_ {i \in \mathrm{J}} \alpha_ {i} | f _ {i} (x) | ^ {2}\right) ^ {1 / 2} \left(\sum_ {i \in \mathrm{J}} \alpha_ {i} | g _ {i} (y) | ^ {2}\right) ^ {1 / 2}\tag{15}
\]

whence also

\[
\left| \sum_ {i \in \mathrm{J}} \alpha_ {i} \overline {{f _ {i} (x)}} g _ {i} (y) \right| \leqslant \mathrm{H} (x, y),\tag{16}
\]

which already establishes property (ii).

By inequality (15), the series \(\sum_{i} \alpha_{i} h_{i}\) converges uniformly on \(\mathrm{K}_{m} \times \mathrm{K}_{n}\) for all \((m, n) \in \mathbf{N}^{2}\). Let \(\widetilde{\mathrm{N}}\) be the function on \(\mathrm{X} \times \mathrm{X}\) defined by

\[
\widetilde {\mathrm{N}} (x, y) = \sum_ {i \in \mathrm{I}} \alpha_ {i} h _ {i} (x, y) = \sum_ {i \in \mathrm{I}} \alpha_ {i} \overline {{f _ {i} (x)}} g _ {i} (y)
\]

for all \((x,y)\in\widetilde{\mathrm{X}}\times\widetilde{\mathrm{X}}\) and by \(\widetilde{\mathrm{N}}(x,y)=0\) otherwise. The function \(\widetilde{N}\) is measurable (INT, IV, p. 175, § 5, n°4, th. 2), and it is continuous on \(K_{m}\times K_{n}\) for all \((m,n)\in\mathbf{N}^{2}\) .

From (16), we have \(|\widetilde{\mathrm{N}}(x,y)| \leqslant \mathrm{H}(x,y)\) for all \((x,y) \in \mathrm{X} \times \mathrm{X}\) , and in particular \(\widetilde{\mathrm{N}}\) belongs to \(\mathcal{L}^2(\mathrm{X} \times \mathrm{X})\) (INT, IV, p. 84, § 5, no. 6, th. 5).

Let us prove that \(N = \widetilde{N}\) on \(\widetilde{X} \times \widetilde{X}\) , which will establish property (i). Let f and g be elements of \(\mathcal{L}^{2}(X)\) . We have

\[
\langle g \mid u _ {\widetilde {N}} (f) \rangle = \int_ {X \times X} \widetilde {N} (x, y) f (x) \overline {{g (y)}} d (\mu \otimes \mu) (x, y)
\]

(formula (12) of IV, p. 172). For every \((x,y)\in \widetilde{\mathrm{X}}\times \widetilde{\mathrm{X}}\) and every finite subset J of I, we have

\[
\left| \sum_ {i \in J} \alpha_ {i} \overline {{f _ {i} (x)}} g _ {i} (y) f (x) \overline {{g (y)}} \right| \leqslant | f (x) g (y) | H (x, y)
\]

by formula (16). Since the right-hand side of this inequality is integrable on X × X (INT, V, p. 95, § 8, n° 3, cor. 2), one can apply Lebesgue's theorem (INT, IV, p. 137, § 3, n° 7, th. 6) and formula (13) to deduce that

\[
\begin{array}{c} \langle g | u _ {\widetilde {N}} (f) \rangle = \sum_ {i \in I} \alpha_ {i} \int_ {X \times X} \overline {{f _ {i} (x)}} g _ {i} (y) f (x) \overline {{g (y)}} d (\mu \otimes \mu) (x, y) \\ = \sum_ {i \in I} \alpha_ {i} \langle f _ {i} | f \rangle \langle g | g _ {i} \rangle = \langle g | u _ {N} (f) \rangle . \end{array}
\]

Therefore, it follows that \(u_{\widetilde{\mathrm{N}}} = u_{\mathrm{N}}\) , whence \(\mathrm{N} = \widetilde{\mathrm{N}}\) in \(\mathrm{L}^2(\mathrm{X} \times \mathrm{X})\) (prop. 3 of III, p. 28, b)).

For every \((m,n)\in\mathbf{N}^{2}\), the functions N and \(\widetilde{N}\) are continuous on \(K_{m}\times K_{n}\), therefore are equal on \(K_{m}\times K_{n}\) (prop. 9 of INT, III, p. 69, § 2, n°2) since the support of the measure induced by \(\mu\otimes\mu\) onto \(K_{m}\times K_{n}\) is equal to \(K_{m}\times K_{n}\). Therefore N coincides with \(\widetilde{N}\) on \(\widetilde{X}\times\widetilde{X}\) .

Finally, the upper bound \(\sqrt{\mathrm{FG}} \leqslant (\mathrm{F} + \mathrm{G}) / 2\) implies that the function \(x \mapsto \sqrt{\mathrm{F}(x)\mathrm{G}(x)} = \mathrm{H}(x,x)\) is integrable on X, whence property (iii).

In the remainder of this issue, we keep the notation of the proposition.

Theorem 2. --- Suppose that \(u_{\mathrm{N}}\) has finite trace. Then the function \(x \mapsto \mathrm{N}(x, x)\) belongs to \(\mathcal{L}^1(\mathrm{X})\) and we have

\[
\mathrm{Tr} (u _ {\mathrm{N}}) = \int_ {\mathrm{X}} \mathrm{N} (x, x) d \mu (x).
\]

By conditions (i) and (ii), we have \(|\mathrm{N}(x,x)| \leqslant \mathrm{H}(x,x)\) for \(x \in \widetilde{X}\) , hence the function \(x \mapsto \mathrm{N}(x,x)\) belongs to \(\mathcal{L}^{1}(\mathrm{X})\) by condition (iii).

Condition (i) proves the equality.

\[
\mathrm{N} (x, x) = \sum_ {i \in \mathrm{I}} \alpha_ {i} \overline {{f _ {i} (x)}} g _ {i} (x)
\]

for all \(x \in \widetilde{\mathrm{X}}\) . By conditions (ii) and (iii), one may apply Lebesgue's theorem (INT, IV, p. 137, § 3, no. 7, th. 6), from which it follows that

\[
\begin{array}{r} \int_ {\mathrm{X}} \mathrm{N} (x, x) d \mu (x) = \sum_ {i \in \mathrm{I}} \alpha_ {i} \int_ {\mathrm{X}} \overline {{f _ {i} (x)}} g _ {i} (x) d \mu (x) \\ = \sum_ {i \in \mathrm{I}} \alpha_ {i} \langle f _ {i} | g _ {i} \rangle = \mathrm{Tr} (u _ {\mathrm{N}}), \end{array}
\]

the last equality resulting from the corollary of Proposition 17 of IV, p. 170.

Lemma 11. --- If the endomorphism \(u_{\mathrm{N}}\) of \(\mathrm{L}^2(\mathrm{X})\) is positive, then \(\mathrm{N}(x,x) \geqslant 0\) for all \(x\) in the support of \(\mu\) .

Since \(u_{N}\) is positive, the families \((f_{i})\) and \((g_{i})\) can be chosen such that \(f_{i} = g_{i}\) for all \(i \in I\) (for remark 3 of IV, p. 151). For every \(x \in \widetilde{X}\) , we then have

\[
\mathrm{N} (x, x) = \sum_ {i \in \mathrm{I}} \alpha_ {i} | f _ {i} (x) | ^ {2} \geqslant 0.
\]

Since \(N\) is continuous and \(\mu(\mathbf{X}-\widetilde{\mathbf{X}})=0\), the function \(x \mapsto N(x,x)\) is nonnegative on the support of \(\mu\) .

Remark. --- The converse assertion of the lemma is not valid (exercise 14 of IV, p. 316).

PROPOSITION 22. --- Suppose that \(u_{\mathrm{N}}\) is a positive endomorphism of \(\mathrm{L}^2(\mathrm{X})\). Then the endomorphism \(u_{\mathrm{N}}\) has finite trace if and only if the function \(x \mapsto \mathrm{N}(x,x)\) is \(\mu\) -integrable. In this case, we have

\[
\mathrm{Tr} (u _ {\mathrm{N}}) = \int_ {\mathrm{X}} \mathrm{N} (x, x) d \mu (x).
\]

If \(u_{N}\) has finite trace, Theorem 2 implies that \(x \mapsto \mathrm{N}(x, x)\) is integrable on X and that its integral is the trace of \(u_{N}\) .

Conversely, suppose that the function \(x \mapsto \mathrm{N}(x, x)\) is integrable. Since \(u_{N}\) is positive, the orthonormal families \((f_i)_{i \in I}\) and \((g_i)_{i \in I}\) can be chosen such that \(f_i = g_i\) for all \(i \in I\) (remark 3 of IV, p. 151). For every \(x \in \widetilde{X}\) , we have

\[
\mathrm{N} (x, x) = \sum_ {i \in \mathrm{I}} \alpha_ {i} | f _ {i} (x) | ^ {2},
\]

whence, for every finite subset J of I, we deduce

\[
\begin{array}{c} \sum_ {i \in \mathrm{J}} \alpha_ {i} = \sum_ {i \in \mathrm{J}} \alpha_ {i} \int_ {\mathrm{X}} | f _ {i} (x) | ^ {2} d \mu (x) \\ = \int_ {\mathrm{X}} \sum_ {i \in \mathrm{J}} \alpha_ {i} | f _ {i} (x) | ^ {2} d \mu (x) \leqslant \int_ {\mathrm{X}} \mathrm{N} (x, x) d \mu (x). \end{array}
\]

The family \((\alpha_{i})_{i\in\mathrm{I}}\) is therefore summable, which implies that the endomorphism \(u_{N}\) has finite trace (prop. 12 of IV, p. 164).

Remark. --- Even when X is compact and N is continuous, the endomorphism \(u_{\mathrm{N}}\) of \(\mathrm{L}^{2}(\mathrm{X})\) is not always of finite trace (cf. exercise 8 of IV, p. 314).

\section{NORMAL ENDOMORPHISMS}

Throughout this paragraph, the Hilbert spaces considered are complex, unless stated otherwise.

\subsection{\texorpdfstring{Complements on spaces \(L^{p}(X,\mu)\)}{Further results on L spaces\^{}\{p\}(X,\textbackslash mu)}}

PROPOSITION 1. --- Let X be a locally compact topological space. Let \(\mu\) be a positive measure on X and \(p \in [1, +\infty[\). Let \((\mathrm{X}_i)_{i \in \mathrm{I}}\) be a locally countable family (INT, IV, p. 190, § 5, no. 9, def. 7) of pairwise disjoint locally compact subsets of X such that the complement of the union of the \(\mathrm{X}_i\) is locally \(\mu\)-negligible. Let \(\varphi_i\) be the characteristic function of \(\mathrm{X}_i\) and \(\mu_i\) be the measure induced by \(\mu\) on \(\mathrm{X}_i\) (INT, IV, p. 186, §5, n° 7, def. 4).

\begin{enumerate}
\def\labelenumi{\alph{enumi})}
\item
  The map \(p_{i}\) : \(f \mapsto f\varphi_{i}\) is a projector of \(\mathcal{L}^{p}(\mathrm{X}, \mu)\) and, by passage to quotients, defines a projector \(\widetilde{p}_{i}\) of \(\mathrm{L}^{p}(\mathrm{X}, \mu)\) . The latter is an orthogonal projector of \(\mathrm{L}^{2}(\mathrm{X}, \mu)\) if p = 2;
\item
  The map \(f \mapsto f|X_i\) induces an isometric isomorphism from the image of \(p_i\) onto the space \(\mathcal{L}^p(X_i, \mu_i)\) and, by passage to quotients, defines an isometric isomorphism from the image of \(\widetilde{p}_i\) onto \(\mathrm{L}^p(X_i, \mu_i)\) , by means of which these two spaces are identified;
\item
  The sum of the spaces \(\mathrm{L}^{p}(\mathrm{X}_{i},\mu_{i})\) is total in the space \(\mathrm{L}^{p}(\mathrm{X},\mu)\) . In the case p = 2, the space \(\mathrm{L}^{2}(\mathrm{X},\mu)\) is the Hilbert sum of the spaces \(\mathrm{L}^{2}(\mathrm{X}_{i},\mu_{i})\) .
\end{enumerate}

The assertion \(a\) ) is elementary. The assertion \(b\) ) follows from the scholium of INT, V, p. 84, § 7, n° 1.

Let \(q \in ]1, +\infty]\) be the conjugate exponent of \(p\) . One identifies the dual of \(L^p(X, \mu)\) with \(L^q(X, \mu)\) (INT, V, p. 61, § 5, n° 8, th. 4). Let \(f\) be a function in \(\mathcal{L}^q(X, \mu)\) whose class \(\widetilde{f}\) in \(L^q(X, \mu)\) satisfies \(\langle \widetilde{f}, \widetilde{p}_i(\varphi) \rangle = 0\) for all \(i \in I\) and every \(\varphi \in L^p(X, \mu)\) . Since the image of \(p_i\) contains \(\mathcal{K}(X_i)\) , it follows that the measure \((f|X_i) \cdot \mu_i\) on \(X_i\) is zero, which means that the restriction of \(f\) to \(X_i\) is locally \(\mu_i\) -negligible on \(X_i\) (INT, V, p. 46, § 5, n° 3, cor. 2). Since the complement in X of the union of the \(X_i\) is locally \(\mu\) -negligible, one concludes that the function \(f\) is locally \(\mu\) -negligible on X (INT, IV, p. 190, § 5, n° 9 and p. 172, n° 2, prop. 5). The class of \(f\) is then zero in \(L^q(X, \mu)\) (using INT, V, p. 8, § 1, n° 3, lemma 1 and the corollary of proposition 9 when \(p \neq 1\) ). By the Hahn--Banach theorem (EVT, II, p. 46, cor. 1), the first part of assertion \(c\) ) follows.

If \(p = 2\) , the image of \(p_i\) is orthogonal to that of \(p_j\) for all \(i \neq j\) since \(X_i\) and \(X_j\) are then disjoint, whence the last assertion.

PROPOSITION 2. --- Let X be a locally compact space that is countable at infinity. Let \(\mu\) be a positive measure on X. For every \(p \in [1, +\infty[\), the space \(\mathrm{L}^{p}(\mathrm{X},\mu)\) is of countable type.

Let \((\mathrm{U}_{n})_{n\in\mathbf{N}}\) be a sequence of relatively compact open sets of X whose union is equal to X and which satisfy \(\overline{U}_{n}\subset U_{n+1}\) for all \(n\in\mathbf{N}\) (TG, I, p. 68, prop. 15). For every \(n\in\mathbf{N}\), the space \(\mathcal{K}(\mathrm{X},\overline{\mathrm{U}}_{n})\) is identified with a closed subspace of the Banach space \(\mathcal{C}(\overline{\mathrm{U}}_{n})\) (INT, III, p. 40, § 1, n° 1); since the latter space is of countable type (TG, X, p. 25, corollary), the same is true of \(\mathcal{K}(\mathrm{X},\overline{\mathrm{U}}_{n})\) (TG, IX, p. 19, cor., (i)). Let \(F_{n}\) be a countable dense subset of \(\mathcal{K}(\mathrm{X},\overline{\mathrm{U}}_{n})\) .

Let \(f \in \mathcal{L}^p(\mathrm{X},\mu)\) and let \(\varepsilon > 0\) . There exists an integer \(n \in \mathbf{N}\) such that \(\int_{\mathrm{X}-\mathrm{U}_n}|f|^p < \varepsilon / 2\) , and there exists \(g \in \mathcal{F}_n\) such that \(\int_{\overline{\mathrm{U}}_n}|f - g|^p < \varepsilon / 2\) . The union of the classes in \(\mathrm{L}^p(\mathrm{X},\mu)\) of the elements of the sets \(\mathcal{F}_n\) is therefore dense in \(\mathrm{L}^p(\mathrm{X},\mu)\) , which concludes the proof (TG, IX, p. 18, prop. 12).

\subsection{Essential image of a measurable function}

In this issue, we denote by X a locally compact topological space, and by \(\mu\) a positive measure on X.

DEFINITION 1. --- Let Y be a topological space. For every measurable function g from X to Y, the \(\mu\)-essential image of \(g\) is the subset of the \(y \in Y\) such that, for every open neighborhood \(U\) of \(y\), the set \(g^{-1}(\mathrm{U})\) is not locally \(\mu\)-negligible in X (INT, IV, p. 172, § 5, no. 2, def. 3).

Lemma 1. --- Let \(g\) be a \(\mu\)-measurable function from X into a topological space Y and let \(S\) be its \(\mu\)-essential image. The set of elements \(x \in X\) such that \(g(x)\) does not belong to S is locally \(\mu\)-negligible in X.

Let \(Z = g^{-1}(Y - S)\) be the set in question.

Suppose first that X is compact and g is continuous. In this case, the pushforward measure \(g(\mu)\) is defined (INT, V, p. 69, § 6, no. 1, def. 1) since \(\mu\) is a bounded measure. It follows from the definitions that the \(\mu\)-essential image of \(g\) is the support of the measure \(g(\mu)\) (cf. INT, V, p. 70, § 6, no. 2, cor. 2), whence \(\mu(Z) = g(\mu)(Y - S) = 0\) (INT, IV, p. 118, § 2, no. 2, prop. 5).

Now consider the general case. Let C be a compact subset of X, and let \(\varepsilon > 0\) be a real number. Since \(g\) is measurable, there exists a compact subset \(C_1 \subset C\) such that \(\mu(C - C_1) \leqslant \varepsilon\) and such that \(g|C_1\) is continuous (INT, IV, p. 169, § 5, no. 1, prop. 1). We then have

\[
\mu (\mathrm{Z} \cap \mathrm{C}) \leqslant \mu (\mathrm{C} - \mathrm{C} _ {1}) + \mu (\mathrm{Z} \cap \mathrm{C} _ {1}) \leqslant \varepsilon + \mu (\mathrm{Z} \cap \mathrm{C} _ {1}).
\]

Let \(\mu_{1}\) be the measure induced by \(\mu\) on \(C_{1}\) (INT, IV, p. 186, § 5, no. 7, def. 4). Let \(S_{1}\) be the \(\mu_{1}\)-essential image of \(g|C_{1}\) . We have \(Z\cap C_{1}\subset(g|C_{1})^{-1}(Y-S_{1})\) . By the first case, the set \(Z\cap C_{1}\) is therefore \(\mu_{1}\) -negligible, and consequently \(\mu\) -negligible (INT, IV, p. 187, § 5, no. 7, lemma 2, (i)). The above inequality becomes \(\mu(Z\cap C)\leqslant\varepsilon\) ; since \(\varepsilon\) and C are arbitrary, the set Z is locally \(\mu\) -negligible (INT, IV, p. 172, § 5, no. 2, prop. 5).

Lemma 2. --- Let \(g\) be a continuous function from X into \(\mathbf{C}\). The \(\mu\)-essential image of \(g\) is the closure of the image of the support of \(\mu\) under \(g\).

Let \(Y\) be the support of \(\mu\). If \(z \in \mathbf{C}\) is not adherent to \(g(Y)\), then there exists an open neighborhood \(U\) of \(z\) such that the open set \(g^{-1}(U)\) does not meet Y and is therefore locally \(\mu\)-negligible; this means that \(z\) does not belong to the \(\mu\)-essential image of \(g\) .

Conversely, if \(z \in \mathbf{C}\) is adherent to \(g(Y)\), then for every open neighborhood \(U\) of z, the set \(g^{-1}(U)\) is an open set in X which meets Y. It is therefore not \(\mu\)-negligible, and since it is open, it is not locally \(\mu\)-negligible (INT, IV, p. 172, § 5, n°2, cor. 2). Hence z belongs to the \(\mu\)-essential image of g.

Lemma 3. --- Let \(g\) be a continuous function from X into \(\mathbf{C}\) and let S be its \(\mu\)-essential image. Suppose that S is nonempty and that \(0 \notin S\), and denote by \(\delta\) the distance from 0 to S. Then \(\delta > 0\).

Set \(h(x) = 1 / g(x)\) if \(g(x) \neq 0\) and \(h(x) = 0\) otherwise. The function \(h\) belongs to \(\mathcal{L}^{\infty}(\mathrm{X},\mu)\). Let \(\widetilde{h}\) be the class of \(h\) in \(\mathrm{L}^{\infty}(\mathrm{X},\mu)\) . We then have the formula \(\delta^{-1} = \| \widetilde{h}\|_{\infty}\) .

Let U be an open neighborhood of 0 such that the open set \(Z = g^{-1}(U)\) is locally \(\mu\) -negligible, hence negligible (INT, IV, p. 172, § 5, no. 2, cor. 2). Let Y be the support of \(\mu\) ; we have \(Y \subset X - Z\) . The restriction of the function \(h\) to \(X - Z\) is continuous and bounded, hence \(h \in \mathcal{L}^{\infty}(X, \mu)\) and the norm of \(\widetilde{h}\) in \(\mathcal{L}^{\infty}(X, \mu)\) is equal to the norm of its restriction to \(X - Z\) . Moreover, for every \(\alpha \in R_{+}\) the set of \(x \in X - Z\) such that \(|h(x)| > \alpha\) is open in \(X - Z\) ; it is therefore locally \(\mu\) -negligible if and only if it does not meet Y (INT, IV, loc. cit. and INT, III, p. 66, § 3, no. 2, def. 1). Consequently, we have

\[
\| \widetilde {h} \| _ {\infty} = \sup _ {x \in \mathrm{Y}} \frac {1}{| g (x) |},
\]

whence

\[
\frac {1}{\| \widetilde {h} \| _ {\infty}} = \inf _ {x \in \mathrm{Y}} | g (x) | = \inf _ {\lambda \in g (\mathrm{Y})} | \lambda | = \inf _ {\lambda \in \overline {{g (\mathrm{Y})}}} | \lambda |
\]

which is equal to δ by the preceding lemma.

\subsection{Universally measurable functions}

In this issue, X is a locally compact topological space. We recall (INT, V, p. 28, § 3, n° 4, def. 2) that a map f from X into a topological space Y is universally measurable if it is \(\mu\) -measurable for every positive measure \(\mu\) on X. It suffices for this that f be \(\mu\) -measurable for every positive measure \(\mu\) with compact support on X (INT, V, p. 28, § 4, no. 3, prop. 6).

Lemma 4. --- Let Y and Z be locally compact topological spaces, \(f\colon X\to Y\) and \(g\colon Y\to Z\) universally measurable mappings. The mapping \(g\circ f\colon X\to Z\) is universally measurable.

Let \(\mu\) be a positive measure with compact support on X, and let C be its support. The restriction of \(f\) to C is \((\mu|\mathrm{C})\)-proper. The map \(g\) is measurable with respect to the image measure \((f|\mathrm{C})(\mu|\mathrm{C})\), so the map \((g \circ f)|\mathrm{C} = g \circ (f|\mathrm{C})\) is \((\mu|\mathrm{C})\)-measurable. Consequently, the map \(g \circ f\) is \(\mu\)-measurable. The lemma follows.

We denote by \(\mathcal{L}_{\mathrm{u}}(\mathrm{X})\) the vector space of complex-valued functions that are universally measurable on X and by \(\mathcal{L}_{\mathrm{u}}^{\infty}(\mathrm{X})\) the subspace of functions \(f \in \mathcal{L}_{\mathrm{u}}(\mathrm{X})\) which are bounded on X. For \(f \in \mathcal{L}_{\mathrm{u}}^{\infty}(\mathrm{X})\), we denote by \(\|f\|_{\infty} = \sup_{x \in X} |f(x)|\) .

Proposition 3. --- a) The set \(\mathcal{L}_{\mathrm{u}}(\mathrm{X})\) is an involutive subalgebra of the involutive algebra of functions from X to C;

\begin{enumerate}
\def\labelenumi{\alph{enumi})}
\setcounter{enumi}{1}
\item
  The set \(\mathcal{L}_{\mathrm{u}}^{\infty}(\mathrm{X})\) is a subalgebra of \(\mathcal{L}_{\mathrm{u}}(\mathrm{X})\) and the mapping defined by \(f\mapsto \| f\|_{\infty}\) is a norm on \(\mathcal{L}_{\mathrm{u}}^{\infty}(\mathrm{X})\) for which \(\mathcal{L}_{\mathrm{u}}^{\infty}(\mathrm{X})\) is a unital star algebra;
\item
  The algebra \(\mathcal{L}_{\mathrm{u}}(\mathrm{X})\) contains \(\mathcal{C}(\mathrm{X})\) and \(\mathcal{L}_{\mathrm{u}}^{\infty}(\mathrm{X})\) contains \(\mathcal{C}_b(\mathrm{X})\) ;
\item
  Every Borel function from X to \(\mathbf{C}\) (TG, IX, p. 61, def. 5) belongs to \(\mathcal{L}_{\mathrm{u}}(\mathrm{X})\) ;
\item
  For any sequence \((f_n)_{n\in \mathbf{N}}\) in \(\mathcal{L}_{\mathrm{u}}(\mathrm{X})\) which converges pointwise to a function \(f\) from X into \(\mathbf{C}\), we have \(f\in \mathcal{L}_{\mathrm{u}}(\mathrm{X})\) .
\end{enumerate}

The assertions \(a\) ) and \(c\) ) follow from the definitions (cf. INT, IV, p. 175, § 5, no. 3, cor. 3). Assertion \(d\) ) is a consequence of INT, IV, p. 179, § 5, no. 5, th. 4, since the inverse image under a Borel function of every Borel subset of \(\mathbf{C}\) is a Borel subset of X, which is universally measurable (INT, V, p. 28, § 3, no. 4). Finally, assertion \(e\) ) follows from Egoroff's theorem (INT, IV, p. 175, § 5, no. 4, th. 2).

To prove assertion \(b\) ), it suffices to note that assertion \(e\) ) implies, \(a\) fortiori, that a uniform limit of universally measurable functions is universally measurable.

There may exist universally measurable functions on X that are not Borel. Indeed, if X is metrizable, the characteristic function of a Souslin subset of X is universally measurable (cf. INT, IV, p. 171, § 5, no. 1, cor. 2); however, in every uncountable Polish space, there exists a Souslin subset that is not Borel ("Souslin's theorem"; this follows from TG, IX, p. 120, exerc. 8 and from the fact that for every uncountable Polish space X, there exists a Borel bijection \(\mathrm{X}\to\mathbf{N}^{\mathbf{N}}\) whose inverse is Borel).

Lemma 5. --- Let \(\mu\) be a positive measure on X. Let \(g\) be a \(\mu\)-measurable function from X into \(\mathbf{C}\) and let \(S\) be its \(\mu\)-essential image. For every function \(f \in \mathcal{L}_{\mathrm{u}}(\mathrm{S})\), the function \(h\) from X into \(\mathbf{C}\) such that \(h(x) = 0\) if \(g(x) \notin \mathrm{S}\) and \(h(x) = f(g(x))\) if \(g(x) \in \mathrm{S}\) is \(\mu\)-measurable.

Let \(x \in X\) and \(U\) be an open relatively compact neighborhood of x. The set \(\mathrm{N} = \mathrm{U} - (\mathrm{U} \cap g^{-1}(\mathrm{S}))\) is \((\mu|\mathrm{U})\)-negligible. For every integer \(n \geqslant 1\), let \(V_{n}\) be an open set such that \(N \subset V_{n} \subset U\) and \(\mu(V_{n} - N) < 1/n\) (INT, IV, p. 116, § 1, no. 4, prop. 19).

The restriction \(\widetilde{g}\) of \(g\) to \(\mathrm{U}-\mathrm{V}_n\) is \(\mu|(\mathrm{U}-\mathrm{V}_n)\)-proper and its image is contained in S. Since \(f\) is universally measurable, the map \(x\mapsto f(g(x))\) from \(\mathrm{U}-\mathrm{V}_n\) into \(\mathbf{C}\) is measurable with respect to the measure \(\mu|(\mathrm{U}-\mathrm{V}_n)\) (INT, V, p. 71, § 6, no. 2, prop. 3). The map \(h_n\) from X into \(\mathbf{C}\) such that \(h_n(x) = 0\) if \(x\notin \mathrm{U}-\mathrm{V}_n\) and \(h_n(x) = f(g(x))\) if \(x\in \mathrm{U}-\mathrm{V}_n\) is therefore \(\mu\)-measurable (INT, IV, p. 193, § 5, no. 10, prop. 16).

Since \(h_{n}(x) \to h(x)\) for \(\mu\)-almost every \(x \in U\), the restriction of \(h\) to U is \(\mu\)-measurable (INT, IV, p. 175, § 5, no. 4, th. 2). One concludes that \(h\) is \(\mu\)-measurable by the localization principle (INT, IV, p. 171, § 4, no. 2, prop. 4).

Remark. --- Under the hypotheses of the lemma, let us denote, by abuse of language, \(f \circ g\) the function \(h\) defined in the lemma.

If \(g = g_{1}\) locally \(\mu\)-almost everywhere on X, the functions \(g\) and \(g_{1}\) have the same \(\mu\)-essential image and one has \(f \circ g = f \circ g_{1}\) .

If \(g\) is a \(\mu\)-measurable function defined \(\mu\)-almost everywhere on X, we shall also denote by \(f \circ g\) the function \(f \circ g_{1}\), where \(g_{1}\) is a \(\mu\)-measurable function defined on X equal to g locally \(\mu\)-almost everywhere.

\subsection{\texorpdfstring{The star algebra \(\mathrm{L}^{\infty}(\mathrm{X},\mu)\)}{The star algebra \textbackslash mathrm\{L\}\^{}\{\textbackslash infty\}(\textbackslash mathrm\{X\},\textbackslash mu)}}

Let X be a locally compact topological space and let \(\mu\) be a positive measure on X. We consider the commutative star algebra \(\mathrm{L}^{\infty}(\mathrm{X},\mu)\) (example 4 of I, p. 103).

Lemma 6. --- Let \(g \in \mathcal{L}^{\infty}(\mathrm{X}, \mu)\). The spectrum of the class of \(g\) in \(\mathrm{L}^{\infty}(\mathrm{X}, \mu)\) is equal to the \(\mu\)-essential image of \(g\) .

Denote by \(\widetilde{g}\) the class of \(g\) in \(\mathrm{L}^{\infty}(\mathrm{X},\mu)\) and \(S\) the \(\mu\)-essential image of \(g\). Let \(z\in \mathbf{C} - \mathbf{S}\). By definition, there exists an open neighborhood \(U\) of \(z\) such that \(\mathrm{Y} = g^{-1}(\mathrm{U})\) is locally \(\mu\)-negligible. The function \(h\) defined by \(h(x) = (g(x) - z)^{-1}\) if \(x\notin \mathrm{Y}\), and \(h(x) = 0\) if \(x\in \mathrm{Y}\), belongs to \(\mathcal{L}^{\infty}(\mathrm{X},\mu)\). Its class \(\widetilde{h}\) in \(\mathrm{L}^{\infty}(\mathrm{X},\mu)\) satisfies \((\widetilde{g} -z)\widetilde{h} = 1\) since \((g(x) - z)h(x) = 1\) for all \(x\) not belonging to the locally \(\mu\)-negligible set Y. Therefore \(z\in \mathbf{C} - \mathrm{Sp}(\widetilde{g})\) .

Conversely, let \(z \in \mathbf{C} - \mathrm{Sp}(\widetilde{g})\). Let \(h \in \mathcal{L}^{\infty}(\mathrm{X},\mu)\) be a function whose class is the inverse of \(\widetilde{g} - z\) in \(\mathrm{L}^{\infty}(\mathrm{X},\mu)\) . There exists a real number \(\mathrm{M} > 0\) such that \(|h(x)| \leqslant \mathrm{M}\) locally \(\mu\)-almost everywhere, and moreover we have \((g(x) - z)h(x) = 1\) locally \(\mu\)-almost everywhere. Let U be the open ball with center \(z\) and radius \(\mathrm{M}^{-1}\) in \(\mathbf{C}\); then \(g^{-1}(\mathrm{U})\) is contained in the locally \(\mu\)-negligible set

\[
h ^ { - 1 } ( ] \mathrm{M} , + \infty [ ) \cup \{ x \in \mathrm{X} \mid ( g ( x ) - z ) h ( x ) \neq 1 \} ,
\]

hence \(z \in \mathbf{C} - \mathrm{S}\) We have proved the lemma.

PROPOSITION 4. --- Let \(g \in \mathcal{L}^{\infty}(\mathrm{X},\mu)\) . Denote by \(\widetilde{g}\) the class of \(g\) in \(\mathrm{L}^{\infty}(\mathrm{X},\mu)\) and \(\mathrm{S}\) the spectrum of \(\widetilde{g}\) . The map \(f \mapsto f \circ g\) (remark, p. 184) is the functional calculus mapping from \(\mathcal{C}(\mathrm{S})\) into \(\mathrm{L}^{\infty}(\mathrm{X},\mu)\) associated with \(\widetilde{g}\) .

Let \(f \in \mathcal{C}(\mathrm{S})\) . The function \(h = f \circ g\) is \(\mu\)-measurable (Lemma 5 of IV, p. 184) and bounded. Denote by \(\widetilde{f \circ g}\) its class in \(\mathrm{L}^{\infty}(\mathrm{X},\mu)\) . The map \(f \mapsto \widetilde{f \circ g}\) is a continuous unital involutive algebra homomorphism from \(\mathcal{C}(\mathrm{S})\) into \(\mathrm{L}^{\infty}(\mathrm{X},\mu)\) . Since \(g(x)\) belongs to \(\mathrm{S}\) locally \(\mu\)-almost everywhere (Lemma 1 of IV, p. 181 and Lemma 6), this homomorphism sends the identity map on \(\mathrm{S}\) to the class \(\widetilde{g}\) of the function \(g\) . It therefore coincides with the functional calculus mapping of \(\widetilde{g}\) (Prop. 7 of I, p. 111).